\documentclass{Configuracoes}  % Usa a classe personalizada
\usepackage[backend=biber,style=abnt,language=brazil]{biblatex}
\addbibresource{referencias.bib}
\usepackage[brazilian]{babel}
% Remove os títulos automáticos do sumário e das listas (os títulos em caixa alta já são
% escritos pela classe; sem isto, apareciam duas vezes)
\newcommand{\removerTitulosAutomaticos}{%
    \renewcommand{\contentsname}{}%
    \renewcommand{\listfigurename}{}%
    \renewcommand{\listtablename}{}%
}
\addto\captionsbrazilian{\removerTitulosAutomaticos}
\AtBeginDocument{\removerTitulosAutomaticos}
\usepackage{tabularx}
\usepackage{array}
\usepackage[table,xcdraw]{xcolor}
\usepackage{colortbl}
\usepackage{comment}
\usepackage{placeins}
\usepackage{array}
\usepackage{xltabular}
\usepackage{booktabs}
\usepackage{placeins}
\usepackage{float}
\newcommand{\fonte}[1]{\vspace{2pt}\raggedright\footnotesize\textit{Fonte: #1}}
\usepackage{newfloat}

\DeclareFloatingEnvironment[
    fileext=loq,
    listname={},
    name=Quadro,
    placement=htbp,
    within=none
]{quadro}

\usepackage{caption}

\captionsetup[figure]{
    labelsep=endash
}

\begin{document}

% === CAPA (não conta e não numera)
\pagenumbering{gobble}
\capa
\clearpage

% === FOLHA DE ROSTO (conta mas não numera)
\pagenumbering{roman}
\thispagestyle{empty}
\folhaRosto
\clearpage

% === FICHA CATALOGRÁFICA (não conta nem numera)
%\pagenumbering{gobble}
%\fichaCatalograficaMacro
%\clearpage

% === PRÉ-TEXTUAIS (conta mas não numera)
\pagenumbering{roman}
\pagestyle{empty}

%\errata
%\clearpage
\folhaAprovacao
\clearpage
\dedicatoria
\clearpage
\sectionCentralizada{AGRADECIMENTOS}

    \vspace{-0.3cm}
    
    \novoBlocoTexto{
         Primeiramente a Deus por ter me dado tantas oportunidades até aqui.
    }
    \vspace{-0.2cm}
    \novoBlocoTexto{
         À Coordenação de Aperfeiçoamento de Pessoal de Nível Superior (CAPES), pelo financiamento da pesquisa.
    }
    \vspace{-0.2cm}
    \novoBlocoTexto{
         Ao professor doutor Carlos Silva, pela orientação segura.
    }
    \vspace{-0.2cm}
    \novoBlocoTexto{
        Aos professores do IFMA, por compartilharem seus conhecimentos e ensinamentos.
    }
   \vspace{-0.2cm}
    \novoBlocoTexto{
        À bibliotecária Maria Silva, pelo apoio técnico na normalização do trabalho.
    }
    \vspace{-0.2cm}
    \novoBlocoTexto{
        À minha família, pelo apoio, pelo amor e pelo reforço incondicional ao longo dessa caminhada, na qual não foi nada fácil construir esse conhecimento. 
    }
    
    

    

\clearpage
\Epigrafe
\clearpage
\sectionCentralizada{RESUMO}

\vspace{-0.3cm}

\novoBlocoPreTextual{
    \noindent 
    As dúvidas dos estudantes constituem informações relevantes para a prática docente, porém nem sempre chegam ao conhecimento dos professores, seja pela resistência dos alunos em manifestá-las, seja pela utilização de ferramentas alternativas de apoio à aprendizagem. Nessa perspectiva, este trabalho teve como objetivo desenvolver e avaliar um protótipo web de assistente educacional baseado em Modelos de Linguagem de Grande Escala (LLM, do inglês Large Language Models), e Geração Aumentada por Recuperação (RAG, do inglês Retrieval-Augmented Generation), destinado à classificação e organização de dúvidas discentes conforme os tópicos do plano de ensino do professor, disponibilizando essas informações em um painel para acompanhamento dos conteúdos questionados pelos estudantes. A pesquisa caracteriza-se como aplicada, quantitativa e exploratória, combinando pesquisa bibliográfica e experimental. O protótipo foi desenvolvido utilizando \textit{Next.js}, \textit{Supabase}, o modelo de linguagem \textit{Qwen2.5-7B-Instruct}, executado localmente, e o modelo de \textit{embeddings} \textit{multilingual-e5-base}. A plataforma permite que estudantes enviem dúvidas por meio de um \textit{chat} e recebam respostas fundamentadas nos materiais didáticos, enquanto os professores visualizam, em um painel, a distribuição das dúvidas por tópico, acompanhada de resumos individuais e perguntas originais sem identificação dos estudantes. Para a avaliação experimental, foram utilizadas 1.416 perguntas de Estrutura de Dados, provenientes de um conjunto de dados público e traduzidas para o português, distribuídas em 12 tópicos. A classificação realizada pela plataforma foi comparada aos tópicos de referência e obteve acurácia de 83,33\%, precisão macro de 73,85\%, revocação macro de 84,76\% e macro F1-score de 74,90\%. Os seis tópicos com mais dúvidas foram os mesmos na referência e na plataforma. Os resultados indicam que a combinação de LLM e RAG permite organizar automaticamente as dúvidas dos estudantes conforme os tópicos do plano de ensino do professor, dispensando a categorização manual de cada dúvida. Essa organização evidencia os tópicos que concentram mais dúvidas e oferece ao docente uma visão das dificuldades da turma em cada conteúdo, servindo como apoio ao planejamento das aulas.
    \noindent 
}
    
\vspace{0.5cm}
\newcommand{\palavrasChave}{\noindent 
        Palavras-chave: inteligência artificial na educação; LLM; RAG; classificação de dúvidas discentes; prática docente.}
\begin{flushleft}
    \palavrasChave
\end{flushleft}
        


% Obs.: as palavras-chave devem ser escolhidas com bastante rigor, pois devem representar adequadamente os principais temas abordados pela pesquisa

\clearpage

% === Abstract: usa arábico (para aparecer certo no sumário), mas não exibe o número
\clearpage
\thispagestyle{empty}
\phantomsection
\sectionCentralizada{ABSTRACT}

\vspace{-0.3cm}
\novoBlocoPreTextual{
    \noindent
    Students’ questions constitute valuable information for teaching practice; however, they do not always reach teachers, either because students are reluctant to express them or because they use alternative tools to support their learning. From this perspective, this study aimed to develop and evaluate a web-based educational assistant prototype using Large Language Models (LLMs) and Retrieval-Augmented Generation (RAG) to classify and organize students’ questions according to the topics outlined in the teacher’s course syllabus. The research is characterized as applied, quantitative, and exploratory, combining bibliographic and experimental research. The prototype was developed using Next.js, Supabase, the locally deployed Qwen2.5-7B-Instruct language model, and the multilingual-e5-base embedding model. The platform allows students to submit questions through a chat interface and receive answers grounded in teaching materials, while teachers can view, through a dashboard, the distribution of questions across topics, accompanied by individual summaries and original questions without identifying the students. For the experimental evaluation, 1,416 Data Structures questions from a public dataset, translated into Portuguese and distributed across 12 topics, were used. The classifications produced by the platform were compared with the reference topic labels, achieving an accuracy of 83.33\% and a macro F1-score of 74,90\%. The analysis revealed higher performance for clearly defined topics and discrepancies related to thematic overlap and categorization criteria. It is concluded that the combination of LLMs and RAG enabled the automatic organization of students’ questions by topic, although the results do not constitute a definitive diagnosis of learning difficulties.
    \noindent
   
}

\vspace{0.5cm}
\newcommand{\keywords}{\noindent 
    Keywords: artificial intelligence in education; LLM; RAG; student question classification; teaching practice}
\begin{flushleft}
    \keywords
\end{flushleft}
\clearpage

% === Restaura ocultação até sumário
\pagestyle{empty}

\listaIlustracoes

\clearpage
\listaTabelas
\clearpage
\listaQuadros
\clearpage
\listaAbreviaturas
\clearpage
%\listaSimbolos
%\clearpage

% === SUMÁRIO
% === Sumário com numeração oculta
% === SUMÁRIO
% === SUMÁRIO
\clearpage
\sectionCentralizada{SUMÁRIO}

\vspace{-0.7cm}
\Sumario
\clearpage

% === Agora armazena a página do sumário
\newcounter{paginasPreTextuais}
\setcounter{paginasPreTextuais}{\value{page}}  % Ex: 13

% === Ativa a numeração arábica visível a partir da introdução
\pagenumbering{arabic}
\setcounter{page}{\value{paginasPreTextuais}}  % Ex: 13
\stepcounter{page}  % Agora página 14

\pagestyle{fancy}
\fancyhf{}
\fancyhead[R]{\fontsize{10pt}{10pt}\selectfont\thepage}
\renewcommand{\headrulewidth}{0pt}

\section{{\fontsize{12pt}{14pt}\selectfont INTRODUÇÃO}}
    
\vspace{-0.3cm}
\novoBlocoTexto {
    Ao estudar qualquer conteúdo, é inevitável que os estudantes encontrem ambiguidades ou dificuldades e precisem de auxílio \parencite{ryan2001avoiding}. Nem toda dúvida, porém, é levada ao professor. Dillon (1988, apud \cite{ryan2001avoiding}) chama atenção para a quantidade de perguntas que permanecem silenciadas, destacando que muitas vezes se prefere evitar a verbalização e aparentar compreensão. 
}
\novoBlocoTexto {
    No ensino superior, muitos estudantes deixam de buscar a ajuda de que precisam por considerá-la constrangedora ou uma admissão de fracasso \parencite{karabenick2011understanding}. Entre os principais motivos para evitar perguntas, destacaram-se a incerteza sobre a qualidade da questão formulada e o receio do julgamento dos colegas. 
}
\novoBlocoTexto {
    Os próprios professores percebem esse silêncio. Em entrevistas com professoras do Ensino Fundamental, \textcite{mertins2022valorizacao} registraram que a vergonha, a timidez e o medo de se expor diante da turma inibem as perguntas, e uma das professoras resumiu: "Se tu só pergunta, alguém tem alguma dúvida? [...] Ninguém, né?". Embora as pesquisas apresentadas tenham sido realizadas em diferentes contextos educacionais, seus resultados apontam para a presença de fatores sociais e emocionais que podem dificultar a manifestação de dúvidas, independentemente das características específicas dos conteúdos estudados.
}
\novoBlocoTexto {
    Para o docente, essa lacuna tem consequências diretas. As professoras entrevistadas por \textcite{mertins2022valorizacao} relatam que as perguntas dos estudantes permitem identificar o que eles sabem e o que ainda não compreenderam, contribuindo para o planejamento das aulas. Uma das participantes afirma: "A partir do momento que eles trazem as perguntas, elaboro a aula" \space \parencite[p.~5]{mertins2022valorizacao}. Quando essas dúvidas não são manifestadas, o professor pode identificar tardiamente as diferenças entre o que foi ensinado e o que foi efetivamente aprendido, o que reduz as possibilidades de intervenção pedagógica \parencite{angelo1993classroom}. 
}
\novoBlocoTexto{
    Além disso, nem toda dúvida que deixa de chegar ao professor permanece sem resposta. Muitos estudantes recorrem inicialmente aos colegas, aos livros ou à internet e procuram o docente apenas quando a dificuldade persiste \parencite{nguyenviet2024exploratory}. Com o avanço das ferramentas de inteligência artificial generativa, surgiram novas possibilidades para o esclarecimento dessas dúvidas. Em relato de experiência com turmas iniciantes de programação que utilizaram o ChatGPT e o DeepSeek, uma das professoras observou que a maioria dos estudantes não buscou ajuda, pois as ferramentas esclareciam suas dúvidas de maneira rápida e eficaz, e manifestou preocupação com a possível dependência desses recursos \parencite{maia2025utilizacao}.
}
\novoBlocoTexto {
    Em pesquisa com estudantes de graduação do Reino Unido, a proporção dos que utilizavam inteligência artificial generativa para obter explicações sobre conceitos passou de 36\%, em 2024, para 58\%, em 2025, mantendo-se como a finalidade de uso mais frequente entre os participantes \parencite{freeman2025student}. No Brasil, 70\% dos estudantes do ensino médio já utilizavam inteligência artificial generativa em pesquisas escolares em 2024 \parencite{nicbr2025tic}. Esses dados não apontam, por si só, uma redução na procura pelos professores, mas sugerem que uma parcela crescente das dúvidas passa a ser tratada em canais aos quais o docente não tem acesso. 
}
\novoBlocoTexto{
    Dessa forma, parte das dificuldades dos estudantes pode permanecer desconhecida pelo professor, seja pela resistência em verbalizar dúvidas, seja pelo uso de fontes alternativas para auxílio. Ainda que esses recursos contribuam para a aprendizagem individual, a falta de informações sobre as dificuldades encontradas limita a compreensão do professor acerca das necessidades de sua turma. 
}
\novoBlocoTexto{
    A partir dessa problemática, estabelece-se a seguinte pergunta de pesquisa: Como proporcionar ao docente acesso organizado às dúvidas dos estudantes, com base em seu próprio material didático, de modo a subsidiar o aprimoramento de suas práticas de ensino?
}
\novoBlocoTexto{
    Para responder a essa questão, o presente trabalho tem como objetivo geral desenvolver e avaliar um protótipo web de um assistente educacional baseado em LLM (sigla em inglês para \textit{Large Language Model}, ou Modelos de Linguagem de Grande Escala) e RAG (sigla em inglês para \textit{Retrieval-Augmented Generation}, ou Geração Aumentada por Recuperação) que classifica e organiza as dúvidas de estudantes com base no plano de ensino do próprio docente, disponibilizando essas informações em um painel para acompanhamento dos conteúdos questionados pelos estudantes.
}

\clearpage
\section{{\fontsize{12pt}{14pt}\selectfont OBJETIVOS}}

\vspace{-0.3cm}
\subsection{{\fontsize{12pt}{14pt}\selectfont Objetivo geral}}
\novoBlocoTexto{
    Desenvolver e avaliar um protótipo web de um assistente educacional baseado em LLM e RAG, em que o RAG recupera trechos dos materiais didáticos de professores e a LLM, com base nesses trechos, classifica e organiza as dúvidas dos estudantes nos tópicos do plano de ensino, disponibilizando essas informações em um painel para acompanhamento dos conteúdos questionados pelos estudantes.
}
\subsection{{\fontsize{12pt}{14pt}\selectfont Objetivos específicos}}
\begin{itemize}
    \item Identificar, por meio da literatura, as dificuldades relacionadas à manifestação, recepção e organização das dúvidas discentes no contexto educacional;
    \item Realizar levantamento de requisitos funcionais e não funcionais do projeto;
    \item Modelar as interfaces de usuário;
    \item Desenvolver um protótipo web baseado em LLM e RAG que permita o registro das dúvidas discentes, a geração de respostas e a classificação dos questionamentos de acordo com os tópicos da disciplina;
    \item Implementar a geração de resumos individuais das dúvidas, e um painel docente para a visualização desses resumos organizados por tópicos com exibição da pergunta original dos estudantes e sem vínculo dos estudantes a estes dados;
    \item Avaliar, em ambiente controlado, o desempenho da classificação temática dasdúvidas realizada pela LLM, com apoio dos trechos recuperados pelo RAG, por meio da comparação com os tópicos de referência de um conjunto de dados público de perguntas de estudantes de Estrutura de Dados.
    
\end{itemize}

\clearpage
\section{{\fontsize{12pt}{14pt}\selectfont REFERENCIAL TEÓRICO}}
\vspace{-0.3cm}
\novoBlocoTexto{
    Este capítulo reúne os conceitos que fundamentam este trabalho e que apoiam a compreensão do assistente proposto e das escolhas feitas nos capítulos seguintes. Ele está organizado em seis seções. A seção~\ref{sec:duvidas-discentes} trata das dúvidas discentes, da busca de ajuda acadêmica e do uso dessas dúvidas como subsídio à prática docente. A seção~\ref{sec:chatbots} apresenta os \textit{chatbots} e o seu papel na educação, tanto no apoio ao estudante quanto como fonte de informação para o professor. As seções~\ref{sec:llm} e~\ref{sec:rag} descrevem as LLM e o RAG, tecnologias em que se baseia o assistente proposto. A seção~\ref{sec:classificacao-duvidas-materiais} aborda a classificação de dúvidas a partir do material didático do professor, com as representações vetoriais, a classificação textual por LLM e as métricas de avaliação. Por fim, a seção~\ref{sec:privacidade} discute a privacidade e a proteção dos dados dos estudantes.
}
\subsection{{\fontsize{12pt}{14pt}\selectfont Dúvidas discentes}}
\label{sec:duvidas-discentes}
\novoBlocoTexto{
    As dúvidas fazem parte do processo de aprendizagem discente, podendo surgir à medida que os estudantes entram em contato com novos conhecimentos e buscam compreendê-los.
}
\subsubsection{{\fontsize{12pt}{14pt}\selectfont Busca de ajuda acadêmica}}
\novoBlocoTexto{
    Diante do cenário de dúvidas discentes, a busca de ajuda acadêmica apresenta-se como uma estratégia de autorregulação que contribui para este processo de aprendizagem \parencite{ryan2001avoiding}. Porém, diferentemente de outras estratégias autorregulatórias, a busca de ajuda acadêmica apresenta uma dimensão interpessoal, uma vez que envolve necessariamente a comunicação ou a interação com outras pessoas \parencite{ryan1997should}.
}
\novoBlocoTexto{
    Historicamente, a busca de ajuda foi associada à dependência e considerada uma prática degradante, frequentemente evitada \parencite{nelsonlegall1985help}. Essa percepção também se manifesta no contexto universitário, no qual muitos estudantes deixam de solicitar o auxílio necessário por considerarem essa atitude constrangedora ou um reconhecimento de fracasso \parencite{karabenick2011understanding}. Entre os fatores que dificultam essa iniciativa, destacam-se o receio de julgamentos por parte de professores e colegas \parencite{ryan1997should}, a indisponibilidade de pessoas que possam oferecer auxílio e o tempo de espera por uma resposta \parencite{ryan2001avoiding}.
}
\novoBlocoTexto{
    Além disso, pesquisas sobre a busca de ajuda em ambientes escolares indicam que estudantes com menor desempenho acadêmico podem apresentar maior resistência em solicitar auxílio, mesmo quando necessitam de apoio \parencite{ryan2001avoiding}. Nesse sentido, a ausência de perguntas não significa necessariamente a inexistência de dificuldades de aprendizagem, pois o estudante pode não compreender suficientemente o conteúdo para formular uma dúvida ou pode sentir constrangimento em manifestá-la \parencite{karabenick2011understanding}. Como consequência, as dificuldades enfrentadas pelos estudantes podem continuar desconhecidas pelos docentes, o que os limita na identificação das necessidades de aprendizagem dos alunos.
}
\subsubsection{{\fontsize{12pt}{14pt}\selectfont Subsídios à prática docente}}
\novoBlocoTexto{
    As dúvidas manifestadas pelos estudantes podem representar necessidades individuais de aprendizagem, como também, informações relevantes para a prática docente. \textcite{freire2013pedagogia} ressaltam que o professor também aprende com aqueles a quem ensina, na medida em que as inquietações e os questionamentos dos estudantes o desafiam a reconsiderar seus próprios conhecimentos. Essa perspectiva relaciona-se ao conceito de avaliação formativa, no qual as evidências sobre a aprendizagem são obtidas, interpretadas e utilizadas para orientar decisões sobre os próximos passos do processo de ensino e aprendizagem \parencite{black2009developing}. 
}
\novoBlocoTexto{
    Em um dos casos relatados por \textcite{angelo1993classroom}, uma professora identificou que sua abordagem estava baseada na forma como havia aprendido o conteúdo durante a pós-graduação, sem considerar suficientemente as necessidades de seus estudantes, o que a levou a modificar as aulas seguintes. Os autores reforçam a necessidade de obter informações continuamente sobre a aprendizagem, considerando que a identificação tardia das dificuldades pode limitar as possibilidades de intervenção pedagógica \parencite{angelo1993classroom}.
}
\novoBlocoTexto{
    Entre as estratégias destinadas à identificação de dificuldades discentes, \textcite{angelo1993classroom} apresentam a técnica denominada \textit{muddiest point}, que consiste em solicitar aos estudantes que indiquem os aspectos menos compreendidos de uma aula ou atividade. As respostas permitem identificar conteúdos que demandam esclarecimentos adicionais e orientar a distribuição do tempo dedicado a determinados tópicos. Para isso, os autores propõem o agrupamento das respostas por assuntos relacionados e a identificação da frequência de ocorrência de cada grupo \parencite{angelo1993classroom}. 
}
\novoBlocoTexto{
    Entretanto, a disponibilidade de informações sobre o andamento da aprendizagem discente não assegura, por si só, sua utilização no aperfeiçoamento do ensino. \textcite{wise2019teaching} evidenciam a complexidade do processo de tomada de decisão docente fundamentada em dados, envolvendo sua interpretação, a definição de ações e a avaliação dos resultados obtidos. Dessa forma, as dúvidas discentes podem carregar subsídios importantes para a prática pedagógica, desde que sejam analisadas e contextualizadas adequadamente. Sua organização pode contribuir para a identificação de conteúdos que exigem maior atenção, enquanto a definição das intervenções didáticas permanece sob responsabilidade do professor.
}
\subsection{{\fontsize{12pt}{14pt}\selectfont\textit{Chatbots}}}
\label{sec:chatbots}
\novoBlocoTexto{
    \textit{Chatbots} são sistemas computacionais desenvolvidos para interagir com usuários por meio da linguagem natural. Também são conhecidos como agentes conversacionais, agentes interativos, \textit{bots} inteligentes ou assistentes digitais \parencite{adamopoulou2020chatbots}.}    
\subsubsection{{\fontsize{12pt}{14pt}\selectfont Função dos \textit{chatbots} no contexto educacional}}
\novoBlocoTexto{
    Na perspectiva educacional, os \textit{chatbots} podem oferecer suporte pedagógico de forma escalável, respondendo dúvidas conceituais e operacionais, oferecendo \textit{feedback} formativo e orientando trajetórias de estudo \parencite{guimaraes2026agentes}.
}
\novoBlocoTexto{
    \textcite{oliveira2026agentes} observam que os modelos LLM ampliaram as possibilidades de personalização e de interação em linguagem natural. Em comparação com sistemas conversacionais baseados em regras predefinidas, os agentes fundamentados em LLM apresentam maior flexibilidade na geração de respostas, possibilitando interações mais adaptáveis às necessidades dos usuários. Nesse cenário, os \textit{chatbots} passaram a ser o tipo de agente educacional mais estudado, presentes em 68,96\% dos 58 estudos revisados por \textcite{elnaffar2026teaching}. 
}
\novoBlocoTexto{
    \textcite{kuhail2023interacting} identificam quatro papéis desses agentes junto ao estudante: o de agente de ensino, que faz o papel do professor ao apresentar instruções, dar exemplos, fazer perguntas e oferecer \textit{feedback} imediato; o de agente par, que atua como colega de estudo; o de agente ensinável, que é ensinado pelo próprio aluno; e o de agente motivacional, que acompanha o estudante e incentiva a aprendizagem. Entre esses papéis, o mais frequente é o de agente de ensino, presente em 55,55\% dos \textit{chatbots} analisados pelos autores \parencite{kuhail2023interacting}.
}
\subsubsection{{\fontsize{12pt}{14pt}\selectfont \textit{Chatbots} no processo de aprendizagem discente}}
\novoBlocoTexto{
    Uma das principais contribuições desses assistentes no processo de aprendizagem discente é oferecer apoio no momento em que a dúvida surge. Na revisão de \textcite{elnaffar2026teaching}, os principais benefícios relatados foram o \textit{feedback} personalizado e o acesso a qualquer hora, que apoiam o estudo no próprio ritmo e reduzem a frustração do aluno. 
}
\novoBlocoTexto{
    Outro papel relevante diz respeito à exposição do estudante. Ao conversar com o assistente, o estudante diminui o receio de ser julgado pelas perguntas que faz \parencite{ma2024integrating}. \textcite{liffiton2023codehelp} relatam que alunos que pareciam desconfortáveis, envergonhados ou relutantes em pedir ajuda ao professor ou ao monitor não tiveram receio de perguntar ao assistente. A própria conversa também contribui para a aprendizagem, pois muitas vezes o estudante não conhece o bastante sobre a própria dúvida para formulá-la adequadamente ou mesmo classificá-la, e o diálogo com o tutor pode ajudá-lo a descobrir qual é, de fato, a sua pergunta \parencite{ma2024integrating}.
}
\novoBlocoTexto{
    Quanto aos efeitos sobre a aprendizagem, uma meta-análise recente aponta resultados favoráveis. \textcite{deng2023metaanalysis} encontraram um efeito global de 0,789 em 32 estudos, com efeito maior quando o \textit{chatbot} atuava como assistente de ensino \parencite{deng2023metaanalysis}. Em um experimento conduzido por \textcite{kestin2025ai}, estudantes que utilizaram um tutor baseado em inteligência artificial (IA, do inglês, \textit{Artificial Intelligence}) apresentaram maior aprendizagem em menor tempo, quando comparados àqueles que participaram de atividades de aprendizagem ativa em sala de aula. Entretanto, os autores ressaltam que o tutor não deve substituir as aulas presenciais, mas preparar os estudantes para aproveitá-las melhor \parencite{kestin2025ai}.
}
\novoBlocoTexto{
    É importante ressaltar que o uso desses assistentes como ferramenta educacional também envolve riscos. Entre os principais desafios identificados por \textcite{elnaffar2026teaching} destacam-se a dependência excessiva, acompanhada de aprendizagem superficial, relatada em 65,52\% dos estudos revisados, seguida das respostas incorretas, inconsistentes ou incoerentes (32,76\%). \textcite{deng2023metaanalysis} registraram ainda em sua meta-análise resultados negativos quanto ao pensamento crítico, ao engajamento e à motivação dos alunos. Diante disso, essas ferramentas não devem substituir o professor, mas sim ser utilizadas de forma complementar ao ensino.
}
\subsubsection{{\fontsize{12pt}{14pt}\selectfont Interações discentes com \textit{chatbots} como subsídios à prática docente}}
\novoBlocoTexto{
    Conforme a descrição dos assistentes na literatura, nota-se que eles se dirigem majoritariamente ao estudante. \textcite{adamopoulou2020chatbots} chegam a mencionar que os \textit{chatbots} reúnem informações ao longo da disciplina que ajudam a melhorar a aprendizagem e o ensino, mas esse uso é pouco explorado. \textcite{alves2024chatbots} encontraram 14 estudos com estudantes e apenas 3 com professores, e apontam a falta de estudos centrados na perspectiva docente como uma lacuna da área \parencite{alves2024chatbots}.
}
\novoBlocoTexto{
    \textcite{gamito2019botanalyzer} observa que, ao oferecer aos alunos um ambiente disponível a qualquer hora para tirar dúvidas, o \textit{chatbot} alivia a carga de trabalho do professor, mas faz com que ele perca contato com os estudantes, o que gera uma lacuna entre as dúvidas que eles têm e o conhecimento que o professor tem delas. Segundo o autor, a única forma de retorno disponível ao docente costuma ser o registro bruto das interações, o que dificulta identificar as dúvidas comuns à turma e os assuntos mais difíceis de serem compreendidos \parencite{gamito2019botanalyzer}. O próprio autor, inclusive, desenvolveu uma ferramenta, o \textit{BotAnalyzer}, que visava facilitar ao professor a visualização das dúvidas e dos problemas que os alunos tiveram durante os estudos, para que ele pudesse se preparar com antecedência para as aulas sem precisar ler um extenso arquivo de registros \parencite{gamito2019botanalyzer}. 
}
\novoBlocoTexto{
    Outras ferramentas também dão ao professor acesso às perguntas dos alunos. O \textit{CodeHelp} é um assistente de programação baseado em LLM, com barreiras de proteção (\textit{guardrails}) que o impedem de entregar as soluções prontas, e foi utilizado durante 12 semanas por 52 estudantes de uma disciplina do primeiro ano de computação e ciência de dados \parencite{liffiton2023codehelp}. Pela leitura das consultas feitas ao assistente, o professor obteve uma visão mais direta do raciocínio dos alunos do que a oferecida pelos trabalhos entregues \parencite{liffiton2023codehelp}. O \textit{MerryQuery}, por sua vez, é um assistente educacional baseado em LLM e RAG que responde aos estudantes com base nos materiais da disciplina, indica as fontes de cada resposta e permite ao docente acompanhar o histórico das conversas \parencite{tabarsi2025merryquery}. Nas três ferramentas, porém, as perguntas não chegam ao professor organizadas pelos conteúdos da disciplina: o \textit{BotAnalyzer}, citado no parágrafo anterior, mostra as palavras mais frequentes nas perguntas, em uma nuvem de palavras \parencite{gamito2019botanalyzer}, enquanto o \textit{CodeHelp} e o \textit{MerryQuery} as apresentam como uma lista. É nessa lacuna que se insere o assistente proposto neste trabalho, que mantém a função de apoio ao estudante, mas tem como objetivo principal levar ao professor as dúvidas dos alunos, organizadas segundo os tópicos extraídos de seu próprio plano de ensino.
}
\subsection{{\fontsize{12pt}{14pt}\selectfont Modelos LLM}}
\label{sec:llm}
\novoBlocoTexto{
    Os modelos LLM são modelos de inteligência artificial geralmente fundamentados na arquitetura \textit{Transformer}, caracterizados pela utilização de redes neurais com bilhões de parâmetros. Esses modelos são pré-treinados em grandes volumes de dados textuais, por meio dos quais aprendem padrões e relações linguísticas, podendo posteriormente ser adaptados para diferentes aplicações de processamento de linguagem natural, como geração de textos, respostas a perguntas e classificação de informações \parencite{zhao2026survey}.
}
\novoBlocoTexto{
    Com o GPT-3, modelo de linguagem com 175 bilhões de parâmetros, \textcite{brown2020language} mostraram que uma LLM suficientemente grande pode realizar tarefas novas sem nenhum treinamento adicional, apenas a partir do texto de entrada. No modo \textit{few-shot}, expressão que pode ser traduzida como ``com poucos exemplos'', o modelo recebe, no próprio texto de entrada, a descrição da tarefa e alguns exemplos já resolvidos, como frases em inglês acompanhadas de suas traduções para o francês, seguidos de um novo caso a resolver, sem que seus parâmetros sejam atualizados. No modo \textit{zero-shot}, ou ``sem exemplos'', o modelo recebe apenas a descrição da tarefa em linguagem natural \parencite{brown2020language}. Esse segundo modo é retomado na seção~\ref{sec:classificacao-textual-llm}.
}
\novoBlocoTexto{
    Essa flexibilidade é um dos fatores que explica a variedade de aplicações dos modelos LLM. Na educação, além de apoiarem os estudantes, eles podem ajudar o professor a personalizar o ensino, a elaborar planos de aula e atividades e a avaliar trabalhos, destacando seus pontos fortes e fracos \parencite{kasneci2023chatgpt}.
}
\novoBlocoTexto{
    Contudo, um dos principais riscos associados ao uso de LLM é a ocorrência de alucinações, caracterizadas pela geração de informações incorretas, inconsistentes ou sem fundamentação factual, mesmo quando apresentadas de maneira coerente e convincente. Segundo \textcite{anhhoang2025survey}, esse fenômeno pode decorrer tanto de problemas na formulação das instruções fornecidas ao modelo quanto de limitações inerentes ao seu funcionamento, relacionadas aos dados utilizados no treinamento e ao processo de geração das respostas. Mesmo os criadores do GPT-3 reconhecem que o modelo às vezes perde a coerência em textos longos e contradiz a si mesmo \parencite{brown2020language}.
}
\novoBlocoTexto{
    Com o objetivo de reduzir a alucinação por meio de embasamento da resposta em fontes verificáveis, surgiu o RAG, que incorpora à LLM conhecimento de bases externas \parencite{gao2023retrieval,lewis2020retrieval}. Porém, é importante ressaltar que a utilização dessa estratégia não elimina a possibilidade de alucinações, embora possa reduzir sua ocorrência. Assim, a literatura recomenda a supervisão humana para acompanhar e avaliar os resultados dessas soluções. \textcite{kasneci2023chatgpt} a consideram indispensável para o uso benéfico das LLM na educação.
}
\subsection{{\fontsize{12pt}{14pt}\selectfont RAG}}
\label{sec:rag}
\novoBlocoTexto{
    O RAG se consolidou como forma de reduzir a alucinação e de incorporar conhecimento atualizado e específico de uma área diante da popularização dos modelos LLM \parencite{gao2023retrieval}. 
    \textcite{lewis2020retrieval} propuseram o RAG para enfrentar limitações dos modelos de linguagem que dependem exclusivamente do conhecimento armazenado em seus parâmetros, como a dificuldade de atualizar esse conhecimento, identificar a origem das informações e evitar a geração de conteúdos incorretos. Essa estratégia combina a memória interna da LLM, chamada paramétrica, com uma memória externa, não paramétrica, formada por documentos acessados por um mecanismo de recuperação \parencite{lewis2020retrieval}. 
}
\novoBlocoTexto{
    Com RAG, primeiro um mecanismo de recuperação busca informações relevantes em uma base de conhecimento externa e depois as fornece à LLM para orientar a geração da resposta. Esse processo contribui para a precisão, a atualidade e a transparência das informações, característica especialmente importante na educação, em que as informações precisam ser verificáveis \parencite{li2025retrieval}. Outra vantagem é que, como o conhecimento fica em uma base externa, ele pode ser atualizado apenas substituindo essa base, sem treinar o modelo novamente \parencite{lewis2020retrieval}. Com isso, o RAG também se destaca por se adaptar a conteúdos que mudam com frequência.
}
\novoBlocoTexto{
    Na educação, o RAG vem sendo aplicado sobretudo em sistemas de perguntas e respostas, aprendizagem personalizada e tutoria, com ganhos de precisão, relevância e personalização das respostas \parencite{arthana2026systematic}. Nessas aplicações, a base de conhecimento é formada pelos próprios materiais didáticos, como livros, planos de ensino e listas de exercícios, que são convertidos em texto antes da indexação \parencite{li2025retrieval}. \textcite{nemeth2025exploring} destacam que a principal vantagem do RAG no ensino superior é justamente o controle do professor sobre a bibliografia usada pelo modelo, que passa a responder com conhecimento fundamentado em materiais selecionados pelo docente. Em uma disciplina introdutória de programação, \textcite{jacobs2024leveraging} usaram as transcrições das aulas como base do RAG, com o objetivo de que a LLM evitasse alucinações e empregasse os mesmos termos técnicos usados pelo professor.
}
\novoBlocoTexto{
    Apesar de suas potencialidades, o RAG apresenta limitações que precisam ser consideradas. A recuperação pode trazer trechos irrelevantes ou deixar de trazer informações essenciais, e o modelo ainda pode gerar conteúdo não sustentado pelo contexto recuperado \parencite{gao2023retrieval}. \textcite{swacha2025retrieval} acrescentam o custo computacional, já que cada pergunta exige uma etapa de recuperação e outra de geração. Essas limitações reforçam a necessidade de supervisão humana na utilização de sistemas educacionais baseados em RAG, a fim de identificar e corrigir possíveis inconsistências.
}
\subsection{{\fontsize{12pt}{14pt}\selectfont Classificação de dúvidas baseada em material didático docente}}
\label{sec:classificacao-duvidas-materiais}
\novoBlocoTexto{
    Por meio da organização de dúvidas por classificação baseada em tópicos ou seções do próprio plano de ensino docente, o professor pode identificar áreas em que os alunos mais pedem ajuda. O docente, então, pode ajustar seus métodos e intervenções, e o tipo de pergunta pode revelar lacunas no currículo ou no material didático \parencite{sheese2023patterns}. O obstáculo é o trabalho: categorizar manualmente todas as perguntas de uma disciplina é uma tarefa tediosa. Por isso, para tal tarefa, a literatura recorre a técnicas automáticas, como a classificação, em que cada texto é atribuído a uma categoria definida previamente.
}
\subsubsection{{\fontsize{12pt}{14pt}\selectfont Representações vetoriais e recuperação semântica}}
\novoBlocoTexto{
    Representações vetoriais e recuperação semântica são métodos para realizar comparação computacional entre textos. A forma mais tradicional de representação de vetores é a contagem de palavras, como a Frequência do Termo–Frequência Inversa do Documento (TF-IDF, do inglês \textit{Term Frequency–Inverse Document Frequency}), que combina a frequência de cada termo no texto com a sua raridade no conjunto, de modo que palavras muito comuns, como “é” ou “que”, pesem menos \parencite{gao2021automatically}. Essa representação, conhecida como “saco de palavras”, não captura diretamente as relações semânticas e contextuais entre as palavras \parencite{grootendorst2022bertopic}. Como alternativa a esse método, consideram-se os \textit{embeddings}, representações vetoriais de palavras, frases ou documentos que levam em conta o contexto em que os termos aparecem, gerados por modelos de linguagem como as Representações de Codificador Bidirecional de Transformadores (BERT, do inglês \textit{Bidirectional Encoder Representations from Transformers}) \parencite{grootendorst2022bertopic}. O BERT considera ao mesmo tempo o contexto à esquerda e à direita de cada palavra \parencite{devlin2019bert}. Com variações como o Sentence-BERT, frases e parágrafos inteiros são convertidos em vetores, e textos de sentido semelhante ficam próximos nesse espaço vetorial \parencite{grootendorst2022bertopic}. São esses mesmos \textit{embeddings} que o RAG usa para encontrar, no material do professor, os trechos mais parecidos com a pergunta do estudante.
}
\novoBlocoTexto{
    A proximidade entre dois vetores pode ser medida pela similaridade de cosseno, que compara a direção dos vetores: quanto mais próximo de 1 for o resultado, maior será a similaridade entre as representações, embora isso não garanta que os textos possuam significados equivalentes. \textcite{karpukhin2020dense} definem a semelhança entre a pergunta e cada passagem pelo produto interno dos seus vetores e observam que, para vetores de comprimento unitário, essa medida equivale ao cosseno. Na recuperação densa, apresentada na seção~\ref{sec:rag}, todas as passagens do material são convertidas em vetores e indexadas com antecedência; quando chega uma pergunta, o sistema calcula o vetor dela e devolve as \textit{k} passagens cujos vetores estão mais próximos, sendo k um número pequeno em relação ao total de passagens \parencite{karpukhin2020dense}.
}
\novoBlocoTexto{
    No entanto, a maior parte dos modelos de \textit{embeddings} foi treinada apenas com textos em inglês, o que limita seu uso em outros idiomas \parencite{wang2024multilingual}. Para contornar essa limitação, \textcite{wang2024multilingual} desenvolveram os modelos \textit{Multilingual E5}, pré-treinados com cerca de 1 bilhão de pares de textos em diversos idiomas e depois ajustados com dados rotulados de alta qualidade, e os avaliaram em recuperação de informação em 16 idiomas e em alinhamento de traduções em mais de 100 idiomas. Segundo os autores, esses modelos podem ser usados em tarefas de recuperação de informação e de similaridade semântica em diferentes idiomas \parencite{wang2024multilingual}. 
}
\subsubsection{{\fontsize{12pt}{14pt}\selectfont Classificação textual baseada em LLM}}
\label{sec:classificacao-textual-llm}
\novoBlocoTexto{
    Uma das metodologias aplicadas à classificação de texto é a classificação zero-shot, na qual o modelo classifica um texto sem nenhum exemplo de treino, contando apenas com os rótulos das categorias \parencite{lobo2025democratizing}. Com LLM, isso é feito por meio de um comando (prompt). \textcite{lobo2025democratizing} informaram à LLM seu papel de analista, a descrição da disciplina e a lista de categorias, e pediram que ela devolvesse apenas o número da categoria de cada pergunta. Em um estudo sobre classificação de perguntas educacionais, \textcite{alfaraby2024analysis} utilizaram o ChatGPT para categorizar 4.959 perguntas em dez classes. Os autores verificaram que o melhor resultado veio do prompt que trazia, além da lista de categorias, a descrição de cada uma, e que essa abordagem é especialmente útil quando não há dados de treinamento \parencite{alfaraby2024analysis}.
}
\novoBlocoTexto{
    É nesse ponto que o material estruturado do professor pode fornecer apoio relevante. Se os tópicos da disciplina forem utilizados como lista fechada de categorias, o problema deixa de ser descobrir os assuntos e passa a ser classificar cada dúvida em assuntos que o professor já reconhece. \textcite{qiu2025systematic} adotaram essa estratégia com uma LLM, instruída a rotular cada conversa com o tópico mais relacionado de uma lista de doze tópicos do curso ou com “Outros” quando não houvesse correspondência. A contagem por tópico revelou volume particularmente alto de mensagens em Probabilidade Condicional, o que coincidia com a dificuldade que o professor já percebia em sala \parencite{qiu2025systematic}. O material oferece ainda a descrição de cada tópico, que melhora a classificação quando incluída no prompt \parencite{alfaraby2024analysis}, e, por meio do RAG, os trechos recuperados carregam metadados como capítulo e página \parencite{gao2023retrieval}, que indicam de qual parte do material a dúvida se aproxima. Para que a LLM não invente um tópico inexistente, a classificação pode ser restringida para que a saída siga obrigatoriamente uma estrutura definida \parencite{geng2023grammar}, como a lista de tópicos do plano de ensino docente.
}
\novoBlocoTexto{
    O desempenho da classificação pode ser avaliado por meio da comparação entre os rótulos atribuídos pelo sistema e aqueles estabelecidos em um conjunto de referência. Entre as métricas utilizadas, destaca-se a acurácia, que representa a proporção total de classificações corretas. A precisão corresponde à proporção de textos atribuídos a determinada categoria que realmente pertencem a ela, enquanto a revocação indica a proporção de textos dessa categoria corretamente identificados pelo sistema \parencite{sokolova2009systematic}. Essas duas medidas são combinadas no F1-score, que expressa seu equilíbrio por meio da média harmônica. Além disso, a matriz de confusão permite identificar quais categorias são classificadas corretamente e quais são confundidas pelo sistema.
}
\novoBlocoTexto{
    Contudo, os resultados apresentados anteriormente nesta seção não significam que a classificação feita por LLM seja perfeita. Em 25 tarefas de processamento de linguagem, o \textit{ChatGPT} apresentou, em média, cerca de 25\% menos qualidade que as melhores soluções especializadas, com perdas maiores nas tarefas mais difíceis \parencite{kocon2023chatgpt}. \textcite{sun2023text} identificaram, nos cenários avaliados, desempenho inferior dos modelos utilizados apenas com instruções em comparação a classificadores treinados especificamente para a tarefa, e seu desempenho costuma ser insuficiente para substituir totalmente a anotação humana \parencite{ziems2024can}. Portanto, a classificação automática das dúvidas deve ser apresentada ao professor como subsídio à análise pedagógica, e não como diagnóstico definitivo de dificuldades de aprendizagem.
}
\subsection{{\fontsize{12pt}{14pt}\selectfont Privacidade e proteção de dados}}
\label{sec:privacidade}
\novoBlocoTexto{
    A Lei Geral de Proteção de Dados Pessoais (LGPD), Lei nº 13.709/2018, define dado pessoal como toda ``informação relacionada a pessoa natural identificada ou identificável" \parencite[art.~5º, I]{brasil2018lgpd}. A adoção de sistemas inteligentes no ensino intensifica a coleta e o tratamento de dados pessoais de estudantes \parencite{miranda2026lgpd}, e os registros de interação em plataformas, assim como os dados de desempenho e comportamento, podem revelar aspectos sensíveis da vida escolar \parencite{vidal2026inteligencia}. 
}
\novoBlocoTexto{
    A LGPD ainda limita o tratamento ``ao mínimo necessário para a realização de suas finalidades" \space \parencite[art.~6º, III]{brasil2018lgpd}. Plataformas educacionais que capturam dados além do necessário para a finalidade pedagógica afrontam esse princípio e ampliam o risco de exposição \parencite{miranda2026lgpd}. Para resolução do problema desta pesquisa, o que mais interessa é saber quais partes do material geram mais dúvidas, e não necessariamente quem as enviou. Essa finalidade pode ser cumprida sem que a identidade do estudante seja exposta. Portanto, a apresentação das informações ao professor é realizada sem identificadores diretos dos estudantes.
}
\novoBlocoTexto{
    A análise do destino para onde as dúvidas discentes são enviadas também deve ser considerada importante. \textcite{vidal2026inteligencia} recomendam verificar, antes de adotar uma ferramenta, o local de armazenamento, o compartilhamento com terceiros e a possibilidade de uso dos dados para treinar modelos. Neste trabalho, o modelo de linguagem é executado localmente, evitando o envio das dúvidas a serviços externos de inferência durante seu processamento. Essa escolha pode reduzir a exposição das informações a terceiros, embora não elimine os riscos associados ao tratamento de dados pessoais, especialmente porque o armazenamento dos registros e dos materiais didáticos utiliza serviços externos. Dessa forma, a execução local da LLM constitui uma medida complementar de proteção, sem dispensar mecanismos de controle de acesso e segurança.
}
\novoBlocoTexto{
    Por fim, nenhuma dessas medidas substitui a transparência. As instituições devem informar aos estudantes para que seus dados serão usados, quem terá acesso a eles e como sua identidade será protegida \parencite{slade2013learning}. No assistente proposto, isso significa informar ao estudante que, ao utilizar a plataforma para fazer uma pergunta, o professor verá um resumo da dúvida e a pergunta original, ambos sem a sua identificação, classificados por tópico da disciplina, com a finalidade de gerar subsídios que possam contribuir para a prática docente. Em síntese, a ocultação da identidade dos estudantes desempenha um papel relevante nesta pesquisa para favorecer a manifestação de dúvidas acadêmicas. Entretanto, a apresentação de informações desidentificadas não elimina completamente o risco de identificação indireta, especialmente em turmas pequenas ou diante de questionamentos muito específicos.
}
\clearpage
\section{{\fontsize{12pt}{14pt}\selectfont MÉTODO DA PESQUISA}}
\label{sec:metodo}
\vspace{-0.3cm}
\novoBlocoTexto{
    A pesquisa classifica-se, quanto à abordagem, como quantitativa, uma vez que considera a compreensão dos elementos por meio de ferramentas estatísticas. Quanto à natureza, classifica-se como aplicada, uma vez que busca gerar conhecimento voltado à solução de um problema prático do ensino. Quanto aos objetivos, a pesquisa é exploratória, visando testar uma solução prática. Quanto aos procedimentos técnicos, combina pesquisa bibliográfica e pesquisa experimental. A disciplina de Estrutura de Dados é adotada como estudo de caso, sem que a proposta se restrinja a ela.
}
\novoBlocoTexto{
    Inicialmente, realizou-se um levantamento bibliográfico em artigos, livros e documentos oficiais sobre a busca de ajuda por estudantes, o uso das dúvidas como subsídio à prática docente, assistentes educacionais, LLM, RAG, classificação de textos e proteção de dados na educação. As buscas foram feitas nas bases Scopus, Google Acadêmico, SBC OpenLib, IEEE Xplore, ACM Digital Library e arXiv. 
}
\novoBlocoTexto{
    Na elaboração deste trabalho, utilizou-se o assistente de inteligência artificial generativa Claude, da Anthropic, como apoio nas seguintes atividades: revisão de partes do texto; conferência das citações e das referências; verificação dos cálculos das métricas; montagem dos slides e do plano de ensino usados no experimento, a partir do texto das apostilas citadas (seção~\ref{sec:preparacao-disciplina}); e escrita do script de tradução do conjunto de dados. Todo o conteúdo produzido com esse apoio foi revisado pela autora, que assume integral responsabilidade pelo texto, pelas análises e pelas conclusões, em consonância com a supervisão humana recomendada na seção~\ref{sec:llm}.
}
\novoBlocoTexto{
    O protótipo foi desenvolvido como uma aplicação web com o \textit{framework Next.js} e a linguagem \textit{TypeScript} e o seu código-fonte está disponível em: \url{https://github.com/mayara-gregorio/Projeto-TCC/tree/v1.0-tcc}. O \textit{Supabase} foi utilizado para o banco de dados PostgreSQL, a autenticação dos usuários e o controle de acesso aos dados, de modo que cada aluno visualize apenas as próprias conversas e o professor visualize os resumos das dúvidas de suas disciplinas, como também, as perguntas originais. O professor cadastra a disciplina e convida os alunos mediante um código de acesso; o aluno, após ingressar, envia suas dúvidas em um chat.
}
\novoBlocoTexto{
    Na preparação da plataforma, o professor envia o plano de ensino da disciplina, do qual a plataforma extrai a lista de tópicos por meio da LLM. O professor anexa o material didático associando-o ao tópico correspondente, em formato PDF. O texto desses materiais é extraído e dividido em fragmentos com cerca de 1.000 caracteres, tamanho também adotado por \textcite{qiu2025systematic}, em busca de equilíbrio entre fragmentos maiores, que preservam mais contexto, mas trazem mais ruído, e fragmentos menores, que têm menos ruído, mas podem não transmitir todo o contexto necessário \parencite{gao2023retrieval}. Cada fragmento, identificado pelo tópico a que pertence, é convertido em vetor pelo modelo de \textit{embeddings} \texttt{multilingual-e5-base} e armazenado no próprio banco de dados da aplicação, por meio da extensão \textit{pgvector} do \textit{PostgreSQL} no \textit{Supabase}.
}
\novoBlocoTexto{
    Quando o aluno envia uma dúvida pelo chat, a pergunta é convertida em vetor pelo mesmo modelo, e o mecanismo de recuperação seleciona os cinco fragmentos mais semelhantes a ela (k = 5). Os fragmentos recuperados são entregues à LLM junto com a lista de tópicos da disciplina, e a LLM, com base neles, realiza três tarefas: classifica a dúvida em um dos tópicos, com um grau de confiança entre 0 e 1; responde ao aluno, indicando o material e o slide de onde veio a informação; e produz um título e um resumo individual da dúvida, em terceira pessoa. Se nenhum fragmento for recuperado, a LLM responde com conhecimento geral e avisa o aluno de que a explicação não veio do material do professor. O banco de dados armazena a pergunta original, o tópico atribuído e o resumo. Por fim, o painel do professor apresenta, para cada tópico, a quantidade de dúvidas, os respectivos resumos e as perguntas originais sem identificar os alunos que as enviaram.
}
\novoBlocoTexto{
    A LLM utilizada é o \textit{Qwen2.5-7B-Instruct}, modelo de pesos abertos da família \textit{Qwen2.5}, que reúne versões de 0,5 a 72 bilhões de parâmetros \parencite{cruzcastaneda2025amadeus}, na versão quantizada em 4 bits (q4\_k\_m, formato GGUF), executada localmente por meio do \textit{llama.cpp}. A quantização representa os pesos do modelo com menos bits, o que reduz a memória necessária e acelera a geração, em troca de possível perda de precisão \parencite{chen2026ondevice}, e a execução local evita o envio das dúvidas a serviços externos de inferência. A classificação usa uma lista fechada com os tópicos da disciplina, os mesmos tópicos presentes no conjunto de dados usado na avaliação (seção~\ref{sec:avaliacao-experimental}).
}
\novoBlocoTexto{
    Para a análise dos resultados, adotam-se duas bases: (i) o atendimento aos requisitos funcionais e não funcionais definidos na seção~\ref{sec:levantamento-requisitos}; e (ii) a avaliação experimental descrita na seção~\ref{sec:avaliacao-experimental}, que, em ambiente controlado, compara os tópicos atribuídos pela LLM às perguntas do conjunto de dados público utilizado com os tópicos de referência desse conjunto.
}
\subsection{{\fontsize{12pt}{14pt}\selectfont Levantamento de requisitos}}
\label{sec:levantamento-requisitos}
    
\novoBlocoTexto{
    Os requisitos do protótipo foram definidos a partir do fluxo descrito neste capítulo. Os requisitos funcionais (RF) são apresentados no Quadro \ref{quad:requisitos-funcionais}, e os não funcionais (RNF), no Quadro \ref{quad:requisitos-nao-funcionais}.
}

\begin{quadro}[H]\captionsetup{name=Quadro}
\centering
\caption{Requisitos funcionais}
\label{quad:requisitos-funcionais}

\renewcommand{\arraystretch}{1.3}
\setlength{\tabcolsep}{4pt}

\begin{tabularx}{\textwidth}{
|>{\centering\arraybackslash}m{1.5cm}
|>{\raggedright\arraybackslash}X
|>{\centering\arraybackslash}m{2.6cm}|
}
\hline

\rowcolor[HTML]{C0C0C0}
\textbf{Código} &
\textbf{Requisito} &
\textbf{Classificação} \\ \hline

RF01 & Permitir ao professor cadastrar disciplinas. & Obrigatório \\ \hline

RF02 & Permitir ao professor convidar alunos por código de acesso. & Obrigatório \\ \hline

RF03 & Permitir ao professor enviar o plano de ensino da disciplina, do qual a LLM extrai os tópicos. & Obrigatório \\ \hline

RF04 & Permitir ao professor associar materiais didáticos em PDF a cada tópico da disciplina. & Obrigatório \\ \hline

RF05 & Extrair o texto dos materiais. & Obrigatório \\ \hline

RF06 & Dividir o texto extraído dos materiais didáticos em fragmentos com cerca de 1.000 caracteres. & Obrigatório \\ \hline

RF07 & Gerar os embeddings de cada fragmento extraído dos materiais didáticos. & Obrigatório \\ \hline

RF08 & Permitir ao aluno enviar dúvidas por meio de um chat. & Obrigatório \\ \hline

RF09 & Recuperar os cinco fragmentos do material mais semelhantes a cada dúvida. & Obrigatório \\ \hline

RF10 & Gerar a resposta ao aluno com base nos fragmentos recuperados. & Obrigatório \\ \hline

RF11 & Classificar cada dúvida em um dos tópicos da disciplina. & Obrigatório \\ \hline

RF12 & Gerar um resumo individual de cada dúvida. & Obrigatório \\ \hline

RF13 & Armazenar a pergunta original, o tópico atribuído e o resumo. & Obrigatório \\ \hline

RF14 & Exibir ao professor um painel com a quantidade de dúvidas, os resumos por tópico e a dúvida. & Obrigatório \\ \hline
\end{tabularx}
\vspace{-0.2cm}
{\fontsize{10pt}{12pt}\selectfont
\caption*{Fonte: Elaborado pela autora (2026).} % Legenda de baixo.
}
\end{quadro}

\begin{quadro}[H]\captionsetup{name=Quadro}
\centering
\caption{Requisitos não funcionais}
\label{quad:requisitos-nao-funcionais}

\renewcommand{\arraystretch}{1.3}
\setlength{\tabcolsep}{4pt}

\begin{tabularx}{\textwidth}{
|>{\centering\arraybackslash}m{1.5cm}
|>{\raggedright\arraybackslash}X
|>{\centering\arraybackslash}m{2.6cm}|
}
\hline

\rowcolor[HTML]{C0C0C0}
\textbf{Código} &
\textbf{Requisito} &
\textbf{Prioridade} \\ \hline

RNF01 & O sistema deve impedir que professores acessem os dados de identificação dos alunos associados às dúvidas. & Alta \\ \hline

RNF02 & O sistema deve garantir que cada aluno tenha acesso somente às próprias conversas. & Alta \\ \hline

RNF03 & O sistema deve aplicar políticas de segurança em nível de linha (RLS) no banco de dados Supabase para controlar o acesso aos registros. & Alta \\ \hline

RNF04 & A inferência da LLM deve ocorrer localmente, sem o envio das dúvidas para serviços externos de inferência. & Alta \\ \hline

RNF05 & O modelo de embeddings deve ser capaz de gerar representações vetoriais de textos em português. & Alta \\ \hline

RNF06 & As respostas geradas pela LLM devem ser apresentadas ao aluno em português. & Alta \\ \hline

RNF07 & Os embeddings devem ser armazenados no banco de dados PostgreSQL da aplicação, utilizando a extensão pgvector. & Alta \\ \hline

RNF08 & A aplicação deve ser compatível com navegadores web modernos. & Alta \\ \hline

RNF09 & O sistema deve informar ao aluno, antes do envio da dúvida, que o professor poderá visualizar a dúvida e um resumo da mesma, sem identificação nominal. & Alta \\ \hline
\end{tabularx}
{\fontsize{10pt}{12pt}\selectfont
\caption*{Fonte: Elaborado pela autora (2026).} % Legenda de baixo.
}
\end{quadro}
\FloatBarrier

\subsection{{\fontsize{12pt}{14pt}\selectfont Avaliação experimental}}
\label{sec:avaliacao-experimental}
\novoBlocoTexto{
    A avaliação tem como objetivo comparar a classificação realizada pela plataforma com a classificação manual de um conjunto de dados de referência. Para isso, as perguntas desse conjunto, previamente rotuladas por tópico, são enviadas à própria plataforma, que, por meio de uma LLM com RAG, as classifica, e os tópicos atribuídos pela LLM são comparados aos tópicos atribuídos às mesmas perguntas no conjunto de dados de referência. Avalia-se uma única condição, a classificação com RAG; a qualidade das respostas e dos resumos não faz parte desta avaliação.
}
\subsubsection{{\fontsize{12pt}{14pt}\selectfont Conjunto de dados de referência}}
\label{sec:conjunto-dados-referencia}
\novoBlocoTexto{
    Utiliza-se o conjunto de dados de \textcite{zaman2024dataset}, publicado em acesso aberto no repositório \textit{Mendeley Data}. Ele reúne perguntas de 126 estudantes de graduação da \textit{Independent University}, Bangladesh, registradas de forma anônima em uma aplicação web durante o semestre de verão de 2023, antes e depois das aulas de duas disciplinas de computação \parencite{zaman2024dataset}. Após a remoção de perguntas duplicadas e a correção de gramática e pontuação, as perguntas foram classificadas manualmente pelos autores segundo a disciplina e o tópico \parencite{zaman2024dataset}.
}
\novoBlocoTexto{
    Neste trabalho, utiliza-se apenas o arquivo da disciplina de Estrutura de Dados, com 1.416 perguntas em inglês distribuídas em 12 tópicos (Tabela \ref{tab:distribuicao-de-perguntas-por-topico-no-conjunto-de-referencia}). A distribuição é desbalanceada: o tópico mais frequente reúne 20,8\% das perguntas, e o menos frequente, 1,0\%. Não são utilizados a pontuação segundo a taxonomia de \textit{Bloom}, também presente no arquivo, nem o arquivo de perguntas não relacionadas às disciplinas, já que a plataforma desenvolvida não adota categoria residual.
}
\begin{table}[htbp]
\centering
\caption{Distribuição das perguntas por tópico no conjunto de referência}
\label{tab:distribuicao-de-perguntas-por-topico-no-conjunto-de-referencia}

\renewcommand{\arraystretch}{1.3}
\setlength{\tabcolsep}{4pt}

\begin{tabularx}{\textwidth}{
>{\raggedright\arraybackslash}X
>{\raggedright\arraybackslash}X
>{\centering\arraybackslash}m{2.0cm}
>{\centering\arraybackslash}m{1.0cm}
}
\hline

\multicolumn{1}{c}{\textbf{Tópico no conjunto}} &
\multicolumn{1}{c}{\textbf{Tradução}} &
\multicolumn{1}{c}{\textbf{Perguntas}} &
\multicolumn{1}{c}{\textbf{\%}} \\ \hline

Python Review & Revisão de Python & 295 & 20,8 \\

Binary Search Tree & Árvore binária de busca & 249 & 17,6 \\

Stacks & Pilhas & 196 & 13,8 \\

Hashing & Hashing (Espalhamento) & 163 & 11,5 \\

Recursion & Recursão & 163 & 11,5 \\

Graph & Grafo & 139 & 9,8 \\

Doubly Linked List & Lista duplamente encadeada & 105 & 7,4 \\

Singly Linked List & Lista simplesmente encadeada & 32 & 2,3 \\

Queues & Filas & 24 & 1,7 \\

Lists & Listas & 21 & 1,5 \\

Pointers & Ponteiros & 15 & 1,1 \\

Nodes & Nós & 14 & 1,0 \\ \hline

\textbf{Total} & & \textbf{1.416} & \textbf{100,0} \\ \hline
\end{tabularx}

{\fontsize{10pt}{12pt}\selectfont
\caption*{Fonte: Elaborado pela autora (2026).}
}

\end{table}
\FloatBarrier

\novoBlocoTexto{
    As perguntas foram traduzidas para o português por meio da API do tradutor automático \textit{DeepL} \parencite{deepl2026}, para validação do funcionamento do sistema no idioma referido; cada pergunta traduzida mantém o tópico de referência da original, e eventuais erros de tradução são considerados na análise dos erros. 
}
\novoBlocoTexto{
    Foi observado que alguns dos tópicos do conjunto de dados de referência são próximos entre si, como Nós, Ponteiros e as listas encadeadas, e algumas perguntas dependem do contexto da aula em que foram feitas, como “\textit{How does it differ from a regular doubly linked list?}”. Essas características podem gerar divergências entre a LLM e a referência e são consideradas na análise dos erros.
}   
\subsubsection{{\fontsize{12pt}{14pt}\selectfont Preparação da disciplina na plataforma}}
\label{sec:preparacao-disciplina}
\novoBlocoTexto{
    Para o experimento, cadastrou-se na plataforma uma disciplina designada "Estrutura de Dados", pois o conjunto de dados utilizado no experimento está relacionado a essa disciplina. Como o plano de ensino original da disciplina não está disponível, elaborou-se um plano com os 12 tópicos do conjunto de referência, com os nomes traduzidos da Tabela \ref{tab:distribuicao-de-perguntas-por-topico-no-conjunto-de-referencia}, que foi enviado à plataforma para a extração dos tópicos. Esse plano foi organizado como os planos de ensino dos cursos de graduação do IFMA, com objetivo geral, ementa, objetivos específicos, conteúdo programático, metodologia, recursos e informações relacionadas à avaliação de alunos. Além do título, a plataforma extrai do plano uma breve descrição de cada tópico. 
}
\novoBlocoTexto{
    Como o conjunto de dados utilizado como referência não inclui o material didático da disciplina, elaborou-se, para cada tópico, um conjunto de slides em português a partir de três apostilas: \textcite{levada2021} e \textcite{silva2026}, com exemplos na linguagem Python, para onze tópicos, e \textcite{celes2002apostila}, com exemplos na linguagem C, para o tópico Ponteiros, que não é abordado nas duas primeiras (Tabela \ref{tab:distribuicao-materiais-por-topico}). Os slides reproduzem o texto dessas apostilas, sem acréscimos, com a fonte indicada no rodapé de cada um; apenas a formatação foi adaptada, e as figuras não foram reproduzidas. O material foi elaborado sem consulta às perguntas do conjunto de referência, para que não antecipe as dúvidas a serem classificadas. Os 12 arquivos, em PDF, somam 405 slides.
}

\begin{table}[H]
\centering
\caption{Distribuição dos materiais didáticos por tópico}
\label{tab:distribuicao-materiais-por-topico}

\renewcommand{\arraystretch}{1.3}
\setlength{\tabcolsep}{4pt}

\begin{tabularx}{\textwidth}{
>{\raggedright\arraybackslash}m{5.5cm}
>{\raggedright\arraybackslash}X
>{\centering\arraybackslash}m{2.0cm}
}
\hline

\multicolumn{1}{c}{\textbf{Tópico}} &
\multicolumn{1}{c}{\textbf{Fontes dos slides}} &
\multicolumn{1}{c}{\textbf{Slides}} \\ \hline

Revisão de Python &
\textcite{levada2021}: ``A linguagem Python'', Aulas 1 e 4 &
86 \\
Listas &
\textcite{silva2026}: ``Sequências e listas'' e Laboratório 2 &
13 \\
Ponteiros &
\textcite{celes2002apostila}: seções 4.3, 5.2 e 7.4 &
25 \\
Nós &
\textcite{levada2021}: Aula 6; \textcite{silva2026}: ``Listas encadeadas'' &
7 \\
Lista simplesmente encadeada &
\textcite{levada2021}: Aula 6; \textcite{silva2026}: ``Listas encadeadas'' e Laboratório 5 &
27 \\
Lista duplamente encadeada &
\textcite{levada2021}: Aula 6; \textcite{silva2026}: ``Listas encadeadas'' &
12 \\
Pilhas &
\textcite{levada2021}: Aula 5; \textcite{silva2026}: ``Pilhas e filas'' e Laboratório 3 &
22 \\
Filas &
\textcite{levada2021}: Aula 5; \textcite{silva2026}: ``Pilhas e filas'' e Laboratório 4 &
27 \\
Recursão &
\textcite{levada2021}: Aulas 2 e 3; \textcite{silva2026}: ``Árvores e recursão'' e Laboratório 6 &
17 \\
Hashing (Espalhamento) &
\textcite{silva2026}: ``Tabelas de espalhamento'' e Laboratório 12 &
12 \\
Árvore binária de busca &
\textcite{levada2021}: Aula 7; \textcite{silva2026}: ``Árvores e recursão'' e Laboratório 7 &
56 \\
Grafo &
\textcite{levada2021}: Aulas 8 a 11; \textcite{silva2026}: ``Grafos e redes'' e Laboratórios 9 a 11 &
101 \\ \hline
\textbf{Total} & & \textbf{405} \\ \hline
\end{tabularx}
{\fontsize{10pt}{12pt}\selectfont
\caption*{Fonte: Elaborado pela autora (2026).}
}
\end{table}
\FloatBarrier

\subsubsection{{\fontsize{12pt}{14pt}\selectfont Execução da classificação}}
\label{sec:execucao-classificacao}
\novoBlocoTexto{
    Cada uma das 1.416 perguntas traduzidas é enviada uma única vez à plataforma, pelo chat do aluno, por meio de testes automatizados com a biblioteca \textit{Selenium WebDriver}, que controla o navegador da mesma forma que um usuário \parencite{selenium2026}. Assim, cada pergunta percorre o mesmo fluxo de uma dúvida real: a recuperação dos cinco fragmentos mais semelhantes, a classificação pela LLM, a resposta ao aluno e o resumo da dúvida. A resposta e o resumo são gerados, mas não são avaliados. Para cada pergunta, o banco de dados registra o tópico e o grau de confiança atribuídos pela LLM e os cinco fragmentos recuperados, com a respectiva similaridade, o que permite identificar os tópicos a que esses fragmentos pertencem. Como a classificação considera apenas a pergunta, os tópicos da disciplina e os fragmentos recuperados, sem o histórico da conversa, a ordem de envio das perguntas não interfere no tópico atribuído.
}
\novoBlocoTexto{
    A LLM é configurada com temperatura baixa (0,2), para reduzir a variação entre execuções, já que, mesmo com a mesma configuração, os resultados de uma LLM podem variar a cada execução \parencite{kocon2023chatgpt}. Segundo \textcite{renze2024effect}, variações de temperatura entre 0 e 1 não afetam de forma estatisticamente significativa o desempenho das LLMs na resolução de problemas. Como o arquivo de perguntas não relacionadas às disciplinas não é utilizado (seção~\ref{sec:conjunto-dados-referencia}), a LLM recebe no prompt a instrução de classificar cada dúvida obrigatoriamente em um dos 12 tópicos.
}
\subsubsection{{\fontsize{12pt}{14pt}\selectfont Métricas e análises de erros}}
\label{sec:metricas-analise-erros}
\novoBlocoTexto{
    O desempenho da classificação é medido pela comparação entre o tópico atribuído pela LLM e o tópico de referência de cada pergunta. Calculam-se a acurácia geral e, para cada tópico, a precisão, a revocação e o F1-score, além da matriz de confusão, conforme definidos na seção~\ref{sec:classificacao-textual-llm}. Como os tópicos têm quantidades muito diferentes de perguntas, a precisão, a revocação e o F1-score são também apresentados como médias macro, calculadas pela média simples dos respectivos valores obtidos nos tópicos, atribuindo o mesmo peso a cada um deles \parencite{sokolova2009systematic}. Assim, o F1 macro corresponde à média dos F1-scores dos 12 tópicos. Como parâmetro mínimo de comparação, considera-se a acurácia de 20,8\%, que seria obtida ao atribuir todas as perguntas ao tópico mais frequente.
}
\novoBlocoTexto{
    Por fim, os erros são analisados por tópico. A partir da matriz de confusão, identificam-se os pares de tópicos mais confundidos e examinam-se as perguntas classificadas incorretamente, para verificar se o erro decorre da proximidade entre os tópicos, de perguntas que dependem do contexto da aula, da tradução, da recuperação de fragmentos de outros tópicos ou da própria LLM. Essa análise é necessária porque, no painel, um erro de classificação corresponde a uma dúvida apresentada ao professor no tópico errado.
}
\subsubsection{{\fontsize{12pt}{14pt}\selectfont Aspectos éticos}}
\novoBlocoTexto{
    A avaliação não envolve participantes humanos: utiliza exclusivamente um conjunto de dados público, cujas perguntas foram coletadas de forma anônima e com o consentimento dos participantes \parencite{zaman2024dataset}.
}
\clearpage
\section{{\fontsize{12pt}{14pt}\selectfont RESULTADOS}}
\vspace{-0.3cm}

    
\novoBlocoTexto{
    Os resultados são apresentados em duas partes. A seção~\ref{sec:secao-5-1} apresenta a plataforma desenvolvida e o atendimento aos requisitos levantados na seção~\ref{sec:levantamento-requisitos}, e as seções~\ref{sec:secao-5-2} a~\ref{sec:discussao-limitacoes} apresentam a execução e os resultados da avaliação experimental descrita na seção~\ref{sec:avaliacao-experimental}.
}

\subsection{{\fontsize{12pt}{14pt}\selectfont A plataforma desenvolvida}}
\label{sec:secao-5-1}
\novoBlocoTexto{
    A plataforma foi implementada como uma aplicação web, acessível por meio de navegadores (RNF08). Para representar o processo de utilização da plataforma, foi empregada a notação \textit{Business Process Model and Notation} (BPMN), utilizada para modelar e representar graficamente processos. O fluxo foi organizado em três diagramas BPMN: preparação da disciplina (Figura~\ref{fig:bpmn-preparacao}), envio e processamento das dúvidas (Figura~\ref{fig:bpmn-duvidas}) e consulta ao painel docente (Figura~\ref{fig:bpmn-painel}). Os diagramas apresentam as principais atividades e interações entre professor, aluno e sistema.
}

\bpmnPreparacao

\bpmnDuvidas

\bpmnPainel

\novoBlocoTexto{
    As Figuras de~\ref{fig:tela-acesso} a~\ref{fig:tela-topico-trechos} apresentam suas telas na ordem do fluxo descrito no capítulo~\ref{sec:metodo}: o acesso à plataforma, a disciplina preparada pelo professor, o envio de uma dúvida pelo aluno e a consulta do professor ao painel. O Quadro~\ref{quad:atendimento-requisitos}, ao final desta seção, relaciona cada requisito dos Quadros~\ref{quad:requisitos-funcionais} e~\ref{quad:requisitos-nao-funcionais} à figura que comprova o seu atendimento.
}

\novoBlocoTexto{
    O cadastro e a autenticação são os mesmos para professores e alunos (Figura~\ref{fig:tela-acesso}). 
}
\telaacesso

\novoBlocoTexto{
    O papel de cada usuário é definido na disciplina: o professor a cadastra, e a plataforma gera os códigos de convite com os quais alunos e outros professores ingressam nela (Figura~\ref{fig:tela-disciplina}). 

\teladisciplina
    Ao acessar a página de tópicos e materiais da disciplina (Figura~\ref{fig:tela-topico}) o professor envia o plano de ensino em PDF, a LLM identifica os tópicos e gera uma breve descrição de cada um, que o professor pode editar antes de salvar. Em seguida, o professor associa a cada tópico os materiais em PDF, dos quais a plataforma extrai o texto, divide-o em fragmentos e gera os \textit{embeddings} de cada fragmento. A descrição extraída do plano aparece no início da página de cada tópico, como em Pilhas (Figura~\ref{fig:tela-topico-resumos}), e os arquivos associados aparecem como origem dos slides recuperados (Figura~\ref{fig:tela-topico-trechos}).
}

\telatopico

\novoBlocoTexto{
    O aluno envia suas dúvidas pelo \textit{chat} e recebe a resposta em português (RNF06), elaborada pela LLM com base nos fragmentos recuperados do material e acompanhada da fonte e do slide de onde veio cada informação, como “[Fonte 2, slide 11]” (Figura~\ref{fig:tela-chat}). A classificação e o resumo da dúvida são gerados no mesmo processamento e não aparecem para o aluno. Abaixo da caixa de mensagem, o \textit{chat} informa ao aluno que o professor poderá visualizar resumos das dúvidas enviadas, sem vínculo à identificação do estudante (RNF09), como mostra a Figura~\ref{fig:tela-chat}.
}

\telachat

\novoBlocoTexto{
    Ao acessar o painel, o professor vê o número de disciplinas, o total de dúvidas classificadas, o número de tópicos com dúvidas, a lista dos tópicos com mais dúvidas e um gráfico com a distribuição das dúvidas entre os tópicos. A Figura~\ref{fig:tela-painel} mostra esse painel ao final do experimento, com as 1.416 perguntas classificadas nos 12 tópicos da disciplina.
}

\telapainel

\novoBlocoTexto{
    Ao abrir um tópico, o professor vê a sua descrição, o número de perguntas registradas e, para cada dúvida, um título e um resumo individual, escritos em terceira pessoa, a pergunta original e sem a identificação do aluno (RNF01), como mostra a Figura~\ref{fig:tela-topico-resumos} para o tópico Pilhas. Na mesma página, uma segunda aba apresenta os slides mais recuperados pela busca nas dúvidas do tópico, com o arquivo de origem de cada um (Figura~\ref{fig:tela-topico-trechos}), o que indica ao professor as partes do material mais recuperadas para a resolução das dúvidas.
}

\telatopicoresumos
\telatopicotrechos

\novoBlocoTexto{
    O controle de acesso é aplicado no próprio banco de dados, por políticas de segurança em nível de linha (RLS) que restringem cada conversa e as suas mensagens ao aluno que as criou (RNF02 e RNF03). Além disso, cada rota da aplicação confere se o usuário participa da disciplina, e as instruções enviadas à LLM tratam a dúvida e os fragmentos como dados, com a orientação de ignorar comandos contidos neles.
}

\novoBlocoTexto{
    O Quadro \ref{quad:atendimento-requisitos} resume o atendimento aos requisitos. Todos foram atendidos. O aviso do RNF09 foi incluído na tela do \textit{chat} depois do experimento, o que não altera os resultados, pois as perguntas foram enviadas por automação (seção~\ref{sec:execucao-classificacao}).
}

\begin{quadro}[H]\captionsetup{name=Quadro}
\centering
\caption{Atendimento aos Requisitos}
\label{quad:atendimento-requisitos}

\renewcommand{\arraystretch}{1.3}
\setlength{\tabcolsep}{4pt}

\begin{tabularx}{\textwidth}{
|>{\raggedright\arraybackslash}X
|>{\raggedright\arraybackslash}X|
}
\hline

\rowcolor[HTML]{C0C0C0}
\multicolumn{1}{|c|}{\textbf{Requisito}} &
\multicolumn{1}{c|}{\textbf{Evidência}} \\ \hline

RF01 -- Cadastro de disciplinas &
Figuras~\ref{fig:tela-disciplina} e~\ref{fig:tela-painel} \\ \hline

RF02 -- Convite de alunos por código de acesso &
Figura~\ref{fig:tela-disciplina} \\ \hline

RF03 -- Envio do plano de ensino e extração dos tópicos &
Figura~\ref{fig:tela-topico-resumos} e Seção~\ref{sec:secao-5-2} \\ \hline

RF04 -- Materiais em PDF associados aos tópicos &
Figura~\ref{fig:tela-topico-trechos} e Seção~\ref{sec:secao-5-2} \\ \hline

RF05 -- Extração do texto dos materiais &
Seção~\ref{sec:secao-5-2} \\ \hline

RF06 -- Divisão do texto em fragmentos &
Seção~\ref{sec:secao-5-2} \\ \hline

RF07 -- Geração dos \textit{embeddings} de cada fragmento &
Seção~\ref{sec:secao-5-2} \\ \hline

RF08 -- Envio de dúvidas pelo \textit{chat} &
Figura~\ref{fig:tela-chat} \\ \hline

RF09 -- Recuperação dos cinco fragmentos &
Figuras~\ref{fig:tela-chat} e~\ref{fig:tela-topico-trechos} \\ \hline

RF10 -- Resposta com base nos fragmentos &
Figura~\ref{fig:tela-chat} \\ \hline

RF11 -- Classificação em um dos tópicos &
Figura~\ref{fig:tela-painel} e Seção~\ref{sec:secao-5-3} \\ \hline

RF12 -- Resumo individual de cada dúvida &
Figura~\ref{fig:tela-topico-resumos} \\ \hline

RF13 -- Armazenamento da pergunta, do tópico e do resumo &
Seção~\ref{sec:secao-5-2} \\ \hline

RF14 -- Painel com quantidades e resumos por tópico &
Figuras~\ref{fig:tela-painel} e~\ref{fig:tela-topico-resumos} \\ \hline

RNF01 -- Professor sem acesso aos nomes &
Figura~\ref{fig:tela-topico-resumos} \\ \hline

RNF02 -- Aluno com acesso só às próprias conversas &
Seção~\ref{sec:secao-5-1} \\ \hline

RNF03 -- Controle de acesso por RLS &
Seção~\ref{sec:secao-5-1} \\ \hline

RNF04 -- LLM executada localmente &
Quadro~\ref{quad:configuracao-ambiente-experimental} \\ \hline

RNF05 -- \textit{Embeddings} com suporte ao português &
Quadro~\ref{quad:configuracao-ambiente-experimental} \\ \hline

RNF06 -- Respostas em português &
Figura~\ref{fig:tela-chat} \\ \hline

RNF07 -- Vetores no banco, com \textit{pgvector} &
Quadro~\ref{quad:configuracao-ambiente-experimental} \\ \hline

RNF08 -- Acesso pelo navegador &
Figuras~\ref{fig:tela-acesso},
\ref{fig:tela-disciplina},
\ref{fig:tela-chat},
\ref{fig:tela-painel},
\ref{fig:tela-topico-resumos} e~\ref{fig:tela-topico-trechos} \\ \hline

RNF09 -- Aviso ao aluno sobre o acesso às perguntas &
Figura~\ref{fig:tela-chat}\, \\ \hline

\end{tabularx}
{\fontsize{10pt}{12pt}\selectfont
\caption*{Fonte: Elaborado pela autora (2026).} % Legenda de baixo.
}
\end{quadro}

\FloatBarrier

\subsection{{\fontsize{12pt}{14pt}\selectfont Execução do experimento}}
\label{sec:secao-5-2}
\novoBlocoTexto{
    Na preparação da disciplina, a plataforma extraiu do plano de ensino os 12 tópicos da Tabela~\ref{tab:distribuicao-de-perguntas-por-topico-no-conjunto-de-referencia}, cada um com a sua descrição, e os 12 arquivos de slides da Tabela~\ref{tab:distribuicao-materiais-por-topico}, com 405 slides no total, foram associados aos respectivos tópicos, divididos em fragmentos e convertidos em embeddings (RF03 a RF07).
}

\novoBlocoTexto{
    As 1.416 perguntas traduzidas foram enviadas pelo \textit{chat} a partir de uma única conta de aluno, criada para o teste; por isso, a página de cada tópico registra um único aluno com dúvidas (Figura \ref{fig:tela-topico-resumos}). Todas as perguntas foram classificadas em um dos 12 tópicos da disciplina, e o painel registrou 1.416 dúvidas classificadas, uma para cada pergunta do conjunto (Figura~\ref{fig:tela-painel}). Para cada pergunta, o banco de dados armazenou a pergunta, o tópico atribuído, o grau de confiança, os fragmentos recuperados e o resumo (RF13).
}

\novoBlocoTexto{
    O Quadro \ref{quad:configuracao-ambiente-experimental} apresenta o ambiente de execução. A LLM e o modelo de embeddings foram executados no próprio computador, sem envio das dúvidas a serviços externos de inferência (RNF04), e os vetores foram armazenados no banco de dados da aplicação, com a extensão pgvector (RNF07).
}



\begin{quadro}[H]\captionsetup{name=Quadro}
\centering
\caption{Configuração do ambiente experimental}
\label{quad:configuracao-ambiente-experimental}

\renewcommand{\arraystretch}{1.3}
\setlength{\tabcolsep}{4pt}

\begin{tabularx}{\textwidth}{
|>{\raggedright\arraybackslash}p{3.5cm}
|>{\raggedright\arraybackslash}X|
}
\hline

\rowcolor[HTML]{C0C0C0}
\multicolumn{1}{|c|}{\textbf{Item}} &
\multicolumn{1}{c|}{\textbf{Configuração}} \\ \hline

Computador &
\textit{Notebook} Acer Nitro AN517-54 \\ \hline

Processador &
Intel Core i7-11800H, de 11ª geração, com 8 núcleos e 16 \textit{threads} \\ \hline

Memória RAM &
64 GB \\ \hline

Placa de vídeo &
NVIDIA GeForce RTX 3050 Laptop GPU \\ \hline

Sistema operacional &
Ubuntu 26.04.1 LTS, 64 bits \\ \hline

LLM &
\textit{Qwen2.5-7B-Instruct}, quantizada em \texttt{q4\_k\_m} (formato GGUF), executada pelo \textit{llama.cpp} na placa de vídeo \\ \hline

Temperatura da LLM &
0,2 \\ \hline

Modelo de \textit{embeddings} &
\texttt{intfloat/multilingual-e5-base} (768 dimensões), executado no processador por um serviço local em Python, com \textit{FastAPI} e \textit{sentence-transformers} \\ \hline

Banco de dados &
\textit{Supabase}, em nuvem (\textit{PostgreSQL} com a extensão \textit{pgvector}) \\ \hline

Aplicação &
\textit{Next.js} 16.3.5, \textit{React} 19.2.8 e \textit{TypeScript} 5 \\ \hline

Envio automatizado &
\textit{Selenium WebDriver} 4.52.0, pelo \textit{chat} da plataforma no navegador \\ \hline

\end{tabularx}

{\fontsize{10pt}{12pt}\selectfont
\caption*{Fonte: Elaborado pela autora (2026).} % Legenda de baixo.
}
\end{quadro}

\FloatBarrier

\FloatBarrier

\novoBlocoTexto{
    Ao final da execução, os registros das perguntas foram exportados do banco de dados, com o tópico atribuído pela LLM, e comparados aos tópicos de referência em planilha eletrônica. A planilha com as 1.416 perguntas traduzidas, o tópico de referência e o tópico atribuído pela plataforma está disponível no \textit{Google Drive}\footnote{Disponível em: \url{https://docs.google.com/spreadsheets/d/12FplYIpg0O_vLE_PhnKJ0lyG4CDuXaGG/edit?usp=sharing}.}, serviço de armazenamento de arquivos em nuvem do Google.
}

\subsection{{\fontsize{12pt}{14pt}\selectfont Desempenho da classificação}}
\label{sec:secao-5-3}
\novoBlocoTexto{
    Das 1.416 perguntas do conjunto de referência, a plataforma atribuiu 1.180 ao mesmo tópico da referência, o que corresponde a uma acurácia de 83,33\%. Esse valor é quatro vezes o parâmetro mínimo definido na seção~\ref{sec:metricas-analise-erros}, de 20,8\%, obtido ao atribuir todas as perguntas ao tópico mais frequente. Todas as perguntas foram classificadas em um dos 12 tópicos da disciplina, e a Tabela~\ref{tab:desempenho-classificacao-topico} apresenta o desempenho em cada um deles, as médias macro e a acurácia geral.
}



\begin{table}[htbp]
\centering
\caption{Desempenho da classificação por tópico}
\label{tab:desempenho-classificacao-topico}

\begingroup
\scriptsize
\renewcommand{\arraystretch}{1.3}
\setlength{\tabcolsep}{2pt}

\begin{tabularx}{\textwidth}{
>{\raggedright\arraybackslash}X
>{\centering\arraybackslash}p{2.0cm}
>{\centering\arraybackslash}p{2.5cm}
>{\centering\arraybackslash}p{1.8cm}
>{\centering\arraybackslash}p{1.5cm}
>{\centering\arraybackslash}p{1.8cm}
>{\centering\arraybackslash}p{1.8cm}
}
\hline

\textbf{Tópico} &
\textbf{Perguntas na referência} &
\textbf{Atribuídas pela plataforma} &
\textbf{Acertos} &
\textbf{Precisão (\%)} &
\textbf{Revocação (\%)} &
\textbf{F1-score (\%)} \\ \hline

Revisão de Python & 295 & 137 & 135 & 98,54 & 45,76 & 62,50 \\

Listas & 21 & 85 & 19 & 22,35 & 90,48 & 35,85 \\

Ponteiros & 15 & 36 & 15 & 41,67 & 100,00 & 58,82 \\

Nós & 14 & 8 & 5 & 62,50 & 35,71 & 45,45 \\

Lista simplesmente encadeada & 32 & 47 & 29 & 61,70 & 90,63 & 73,42 \\

Lista duplamente encadeada & 105 & 121 & 102 & 84,30 & 97,14 & 90,27 \\

Pilhas & 196 & 169 & 163 & 96,45 & 83,16 & 89,32 \\

Filas & 24 & 33 & 21 & 63,64 & 87,50 & 73,68 \\

Recursão & 163 & 206 & 160 & 77,67 & 98,16 & 86,72 \\

\textit{Hashing} (Espalhamento) & 163 & 164 & 156 & 95,12 & 95,71 & 95,41 \\

Árvore Binária de Busca & 249 & 262 & 242 & 92,37 & 97,19 & 94,72 \\

Grafo & 139 & 148 & 133 & 89,86 & 95,68 & 92,68 \\ \hline

\textbf{Média macro} &
&
&
&
\textbf{73,85} &
\textbf{84,76} &
\textbf{74,90} \\

\textbf{Total} &
\textbf{1.416} &
\textbf{1.416} &
\textbf{1.180} &
&
&
\\

\textbf{Acurácia geral (\%)} &
\multicolumn{6}{c}{\textbf{83,33} (1.180 acertos em 1.416 perguntas)} \\ \hline

\end{tabularx}
\endgroup
{\fontsize{10pt}{12pt}\selectfont
\caption*{Fonte: Elaborado pela autora (2026).}
}

\end{table}

\FloatBarrier

\novoBlocoTexto{
    O desempenho variou entre os tópicos. \textit{Hashing} (Espalhamento), Árvore binária de busca, Grafo, Lista duplamente encadeada e Pilhas obtiveram F1-score acima de 89\%, enquanto Listas, Nós, Ponteiros e Revisão de Python ficaram entre 35,85\% e 62,50\%. O macro F1-score, calculado pela média simples dos F1-scores de cada tópico e atribuindo o mesmo peso a todos eles (seção~\ref{sec:metricas-analise-erros}), foi de 74,90\%, abaixo da acurácia geral de 83,33\%.
}
\novoBlocoTexto{
    Em dois grupos de tópicos, precisão e revocação seguiram direções opostas. Revisão de Python teve a maior precisão (98,54\%) e uma das menores revocações (45,76\%): quase todas as perguntas atribuídas a esse tópico eram dele, mas 160 das suas 295 perguntas foram atribuídas a outros. Listas e Ponteiros mostraram o inverso, com revocação de 90,48\% e 100\% e precisão de 22,35\% e 41,67\%, porque receberam perguntas de outros tópicos, sobretudo de Revisão de Python. A Figura~\ref{fig:matriz-confusao} apresenta a matriz de confusão, com o destino de cada pergunta.
}

\matrizconfusao

\FloatBarrier

\subsection{{\fontsize{12pt}{14pt}\selectfont Análise dos erros}}
\label{sec:analise-erros}
\novoBlocoTexto{
    Das 236 perguntas atribuídas a um tópico diferente do de referência, 160 (67,8\%) eram de Revisão de Python e 33 (14,0\%) de Pilhas; os outros dez tópicos somaram 43 (18,2\%). Sem as perguntas de Revisão de Python, a concordância nos outros onze tópicos foi de 93,22\% (1.045 de 1.121). Os pares mais frequentes na matriz de confusão foram Revisão de Python → Listas (63 perguntas), Revisão de Python → Recursão (35), Revisão de Python → Ponteiros (11), Revisão de Python → Lista simplesmente encadeada (10) e Pilhas → Recursão (10).
}

\novoBlocoTexto{
    A leitura das perguntas do conjunto aponta três causas prováveis. A primeira é a divisão diferente dos tópicos. No conjunto, Revisão de Python reúne perguntas sobre recursos da linguagem, e 58 das suas 295 perguntas citam listas, vetores (\textit{arrays}), matrizes ou tuplas. No material da disciplina, as listas do Python são o conteúdo do tópico Listas, para o qual a plataforma levou 63 perguntas de Revisão de Python, 39 delas entre essas 58. As outras 24 tratam de outros recursos da linguagem, como laços de repetição, tipos mutáveis e imutáveis, cadeias de caracteres (\textit{strings}) e dicionários, e parte desses assuntos também aparece no material de Listas, que apresenta a lista como sequência mutável e percorre os seus elementos com a função \texttt{enumerate}.
}

\novoBlocoTexto{
    A segunda é que parte dos rótulos parece indicar o tema da aula, e não o assunto da pergunta. Os autores relatam que os estudantes muitas vezes perguntavam sobre assuntos diferentes do tópico escolhido no aplicativo, o que exigiu a categorização manual das perguntas \parencite{zaman2024dataset}. Ainda assim, 44 perguntas rotuladas como Revisão de \textit{Python}, possivelmente, tratam, pelo texto, de outros tópicos: 25 de recursão, como ``Como resolver o problema da Torre de Hanói usando recursão?”, 6 de grafos, 4 de \textit{hashing} e 9 de outros tópicos. A plataforma atribuiu 43 delas ao tópico do assunto, assim como uma pergunta rotulada como Grafo que trata de árvore binária AVL, tipo de árvore binária de busca balanceada cujo nome vem das iniciais dos sobrenomes de seus criadores, Adelson-Velsky e Landis.
}

\novoBlocoTexto{
    A terceira são termos com mais de um sentido. Em Pilhas, 37 das 196 perguntas tratam da pilha como região de memória (\textit{heap}, memória, estouro de pilha) e 19 citam recursão ou a pilha de chamadas. No material, a memória aparece no tópico Ponteiros e a pilha de chamadas no de Recursão, para os quais foram levadas 9 e 10 perguntas de Pilhas; as 9 levadas a Ponteiros tratam todas da pilha como memória do programa, como a alocação de memória na pilha, o estouro de \textit{buffer} e o registrador de ponteiro de pilha. Do mesmo modo, 7 das 14 perguntas de Nós citam árvore, raiz, folha ou filho, e as 6 perguntas de Nós atribuídas a Árvore binária de busca tratam dos nós de uma árvore.
}

\novoBlocoTexto{
    O Quadro \ref{quad:divergencias-classificacao} apresenta um exemplo de pergunta para cada uma dessas causas.
}

\begin{quadro}[H]\captionsetup{name=Quadro}
\centering
\caption{Exemplos de divergências na classificação das perguntas}
\label{quad:divergencias-classificacao}

\renewcommand{\arraystretch}{1.3}
\setlength{\tabcolsep}{3pt}

\begin{tabularx}{\textwidth}{
|>{\raggedright\arraybackslash}p{2.5cm}
|>{\raggedright\arraybackslash}X
|>{\raggedright\arraybackslash}p{2.2cm}
|>{\raggedright\arraybackslash}p{2.0cm}|
}
\hline

\rowcolor[HTML]{C0C0C0}
\textbf{Causa provável} &
\textbf{Pergunta enviada à plataforma} &
\textbf{Tópico de referência} &
\textbf{Tópico atribuído} \\ \hline

Divisão diferente dos tópicos &
Qual é a diferença entre uma lista e uma tupla em Python, tanto em termos de sintaxe quanto de mutabilidade? &
Revisão de Python &
Listas \\ \hline

Rótulo pelo tema da aula &
Como calcular o Máximo Divisor Comum (MDC) de dois números usando recursão? &
Revisão de Python &
Recursão \\ \hline

Termo com mais de um sentido &
Quais são os benefícios e os desafios da alocação de memória na pilha (\textit{stack}) em comparação com a alocação no \textit{heap}? &
Pilhas &
Ponteiros \\ \hline

\end{tabularx}

{\fontsize{10pt}{12pt}\selectfont
\caption*{Fonte: Elaborado pela autora (2026).}
}
\end{quadro}

\FloatBarrier


\FloatBarrier

\novoBlocoTexto{
A etapa em que cada divergência surgiu, na recuperação dos fragmentos ou na escolha da LLM, não foi quantificada nesta avaliação, embora o banco de dados registre os cinco fragmentos recuperados para cada pergunta. O painel mostra, porém, que a busca também recupera fragmentos de outros tópicos: entre os slides mais recuperados nas dúvidas atribuídas a Pilhas, está um slide do material de Filas (Figura \ref{fig:tela-topico-trechos}).
}

\novoBlocoTexto{
Quanto à tradução, a comparação entre o texto original e o traduzido mostrou que os principais termos técnicos foram traduzidos como no material da disciplina em quase todas as perguntas; linked list, por exemplo, foi traduzido como lista encadeada em 148 das 151 perguntas em que aparece. Traduções incomuns, como “lista duplamente articulada” e “percussão” no lugar de percurso (traversal), ocorreram em 23 perguntas, e só 3 delas foram atribuídas a outro tópico. Em duas, a divergência já se explica pelas causas anteriores, e apenas em uma, sobre a detecção de ciclos em uma “lista duplamente articulada”, atribuída a Grafo, a tradução pode ter contribuído. Assim, nos termos conferidos, a tradução não explica as divergências observadas. A ambiguidade do termo pilha, que designa tanto a estrutura de dados quanto a região de memória, já existe no inglês stack e, portanto, não decorre da tradução.
}

\novoBlocoTexto{
Para o professor, essas divergências aparecem na contagem de dúvidas por tópico do painel (Figura \ref{fig:tela-painel}). Os seis tópicos com mais dúvidas são os mesmos na referência e na plataforma (Revisão de Python, Árvore binária de busca, Pilhas, Recursão, \textit{Hashing} (Espalhamento) e Grafo), mas em outra ordem: Revisão de Python, o tópico com mais perguntas na referência (295), aparece no painel em 6º lugar (137), enquanto Recursão passa de 163 para 206 dúvidas e Listas, de 21 para 85. Assim, o painel indica os conteúdos com mais dúvidas, mas subestima Revisão de Python e superestima Listas e Recursão, o que reforça que ele deve ser lido como subsídio, e não como diagnóstico definitivo (seção~\ref{sec:classificacao-textual-llm}).
}
\subsection{{\fontsize{12pt}{14pt}\selectfont Discussão e limitações}}
\label{sec:discussao-limitacoes}
\novoBlocoTexto{
Os resultados indicam que a plataforma organizou as dúvidas pelos tópicos do material do professor com 83,33\% de concordância com a classificação manual de referência. Em relação ao problema desta pesquisa, isso significa que o professor pode reconhecer quais tópicos concentram mais dúvidas: os seis tópicos com mais perguntas foram os mesmos na referência e na plataforma (seção~\ref{sec:analise-erros}). Essa organização dispensa a categorização manual, tarefa trabalhosa que, na prática, impede o professor sem tempo disponível de aproveitar esses dados \parencite{lobo2025democratizing}.
}

\novoBlocoTexto{
O desempenho foi maior nos tópicos de conteúdo bem delimitado, como \textit{Hashing} (Espalhamento), Árvore binária de busca e Grafo, com F1 acima de 92\%, e menor nos tópicos cujo conteúdo se sobrepõe ao de outros, como Revisão de Python, Listas, Nós e Ponteiros. Parte dessas divergências decorre do critério de classificação: a referência agrupa parte das perguntas pelo tema da aula em que foram feitas, enquanto a plataforma as classifica pelo conteúdo do material do professor, que é a proposta deste trabalho. Assim, uma pergunta de Revisão de Python sobre listas atribuída ao tópico Listas conta como erro na avaliação, mas aponta ao professor a parte do material relacionada à dúvida, como ocorreu com 39 das 63 perguntas de Revisão de Python atribuídas a Listas e com 44 das 45 perguntas cujo rótulo difere do assunto (seção~\ref{sec:analise-erros}).
}

\novoBlocoTexto{
Esse resultado deve ser lido considerando a dificuldade da própria tarefa. No estudo de \textcite{qiu2025systematic}, em uma avaliação humana de um subconjunto de 285 itens rotulados por uma LLM, três avaliadores independentes concordaram totalmente em 83,16\% dos itens (237 de 285), com kappa de 0,462, valor que os autores classificam como concordância moderada. Isso sugere que, mesmo entre avaliadores humanos, o julgamento do tópico de uma dúvida não é unânime. \textcite{lobo2025democratizing}, por sua vez, observaram desempenho fraco de modelos com cerca de 3 e 3,8 bilhões de parâmetros e lembram que, no estudo de \textcite{goldberg2021refining}, a concordância entre professores ao rotular perguntas foi de apenas cerca de 65\%; neste trabalho, um modelo de 7 bilhões de parâmetros, quantizado e executado localmente, com apoio do RAG, chegou a 83,33\%. Essas comparações são apenas indicativas, pois as tarefas diferem: em \textcite{qiu2025systematic} e em \textcite{goldberg2021refining}, os valores medem a concordância entre avaliadores, e não entre a LLM e uma referência, e em \textcite{lobo2025democratizing}, classificava-se o tipo da pergunta, e não o tópico.
}

\novoBlocoTexto{
Mesmo com esse nível de concordância, uma em cada seis perguntas (16,67\%) foi apresentada em um tópico diferente do de referência, e as divergências, concentradas em poucos pares de tópicos, alteraram a ordem dos tópicos no painel (seção~\ref{sec:analise-erros}). Por isso, como discutido na seção~\ref{sec:classificacao-textual-llm}, a contagem por tópico deve ser apresentada ao professor como subsídio à análise pedagógica, e não como diagnóstico definitivo, cabendo a ele interpretar os dados e definir as ações, etapas que \textcite{wise2019teaching} apontam como parte da tomada de decisão docente fundamentada em dados.
}

\novoBlocoTexto{
O experimento tem limitações. As perguntas foram escritas em inglês e traduzidas automaticamente, parte dos rótulos de referência indica o tema da aula, e não o assunto da pergunta; os próprios autores reconhecem que o conjunto contém perguntas vagas ou mal estruturadas \parencite{zaman2024dataset}. Os itens do conteúdo programático foram gerados a partir dos tópicos extraídos do conjunto de referência, e as descrições dos tópicos extraídas do plano são enviadas à LLM na classificação, o que pode ter favorecido a concordância obtida (seção~\ref{sec:preparacao-disciplina}). Além disso, avaliou-se uma única disciplina, em uma única execução, com um único modelo e uma única configuração, sem medir a variação entre execuções, e as perguntas foram enviadas por automação, a partir de uma única conta, e não por estudantes de uma turma real. A etapa em que surgiu cada divergência não foi quantificada. Por fim, o material da disciplina foi elaborado a partir de apostilas para o experimento, e não pelo professor da disciplina, e a qualidade das respostas e dos resumos não foi avaliada.
}

\clearpage
\section{{\fontsize{12pt}{14pt}\selectfont CONCLUSÃO}}
\vspace{-0.3cm}
\novoBlocoTexto {
    Este trabalho partiu do problema de pesquisa apresentado na Introdução, formulado na seguinte QP: ``\textit{Como proporcionar ao docente acesso organizado às dúvidas dos estudantes, a partir de seu próprio material didático, de modo a subsidiar o aprimoramento de suas práticas de ensino?}''. Para investigá-la, definiu-se como objetivo geral desenvolver e avaliar um protótipo web de um assistente educacional baseado em LLM e RAG, em o RAG recupera trechos dos materiais didáticos de professores e a LLM, com base nesses trechos, classifica e organiza as dúvidas dos estudantes nos tópicos do plano de ensino, disponibilizando essas informações em um painel para acompanhamento dos conteúdos questionados pelos estudantes.
}
\novoBlocoTexto {
    Os objetivos específicos foram atendidos ao longo do trabalho. O primeiro, identificar na literatura as dificuldades relacionadas à manifestação, recepção e organização das dúvidas discentes, foi cumprido pela revisão da literatura apresentada na Introdução e no referencial teórico, que mostrou que muitas dúvidas não chegam ao professor, seja pela resistência dos estudantes em manifestá-las, seja pelo uso de outras fontes de ajuda, como as ferramentas de inteligência artificial generativa, e que essas dúvidas são subsídios relevantes ao planejamento das aulas. O segundo, o levantamento de requisitos, resultou nos 14 requisitos funcionais e nos 9 requisitos não funcionais da seção~\ref{sec:levantamento-requisitos}. O terceiro, a modelagem das interfaces de usuário, foi atendido com os fluxos BPMN, que organizam as interações entre professor, aluno e sistema, e com as telas apresentadas na seção~\ref{sec:secao-5-1}. O quarto, o desenvolvimento do protótipo, foi realizado com o \textit{framework Next.js}, o \textit{Supabase}, a LLM \textit{Qwen2.5-7B-Instruct}, executada localmente, e o modelo de \textit{embeddings} \textit{multilingual-e5-base}; o protótipo registra as dúvidas, responde aos estudantes com base no material da disciplina e classifica cada dúvida em um dos tópicos do plano de ensino. O quinto foi atendido com a geração dos resumos individuais e com o painel docente organizado por tópicos, que exibe os resumos das perguntas e as perguntas originais sem vínculo com os estudantes. O Quadro~\ref{quad:atendimento-requisitos} mostra o atendimento de todos os requisitos.
}
\novoBlocoTexto {
    O desempenho da classificação foi avaliado em ambiente controlado, atendendo ao sexto objetivo específico deste trabalho. As 1.416 perguntas de estudantes de Estrutura de Dados do conjunto de dados de \textcite{zaman2024dataset}, traduzidas para o português, foram enviadas à plataforma pelo \textit{chat}, e o tópico atribuído pela LLM coincidiu com o tópico de referência em 83,33\% das perguntas, com precisão macro de 73,85\%, revocação macro de 84,76\% e macro F1-score de 74,90\%. O desempenho foi alto nos tópicos de conteúdo bem delimitado, como \textit{Hashing} (Espalhamento), Árvore binária de busca e Grafo, e menor nos tópicos que se sobrepõem, como Revisão de Python, Listas, Nós e Ponteiros. A análise dos erros apontou como causas prováveis das divergências a divisão diferente dos tópicos, os rótulos de referência que indicam o tema da aula, e não o assunto da pergunta, e os termos com mais de um sentido (seção~\ref{sec:analise-erros}); nos termos técnicos examinados, a tradução não se mostrou uma explicação relevante para as divergências. Com o protótipo desenvolvido e avaliado, o objetivo geral foi atendido.
}
\novoBlocoTexto {
    Em resposta à QP, os resultados indicam que é possível proporcionar ao docente acesso organizado às dúvidas dos estudantes, a partir de seu próprio material didático, por meio de um assistente em que o RAG recupera os trechos do material relacionados a cada dúvida e a LLM, com base neles, classifica a dúvida em um dos tópicos do plano de ensino e gera o seu resumo, apresentados ao professor em um painel organizado por tópico e sem a identificação dos estudantes. Essa organização dispensa a categorização manual das perguntas e, no experimento, apontou os mesmos seis tópicos com mais dúvidas da referência, ainda que em outra ordem. Como cerca de uma em cada seis perguntas foi apresentada em um tópico diferente do de referência, o painel deve ser entendido como subsídio à análise pedagógica, e não como diagnóstico definitivo, cabendo ao professor interpretar os dados e decidir as ações que podem aprimorar suas práticas de ensino.
}
\novoBlocoTexto {
    Esses resultados têm limitações: a avaliação usou perguntas traduzidas automaticamente, de uma única disciplina, enviadas por automação a partir de uma única conta, em uma única execução e com um único modelo, e um material elaborado a partir de apostilas para o experimento; além disso, a qualidade das respostas e dos resumos não foi avaliada (seção~\ref{sec:discussao-limitacoes}).
}
\novoBlocoTexto {
    Como trabalhos futuros, propõe-se o agrupamento, por similaridade, das dúvidas de um mesmo tópico, para que o professor visualize de forma reunida perguntas semelhantes feitas por diferentes estudantes, à semelhança do agrupamento por assuntos proposto por \textcite{angelo1993classroom} para as respostas da técnica \textit{muddiest point}. Esse recurso não foi incluído nesta versão porque o conjunto de dados utilizado na avaliação não dispõe de agrupamentos de referência. Também se propõem a inclusão de uma categoria para dúvidas não relacionadas à disciplina, que poderia ser avaliada com o arquivo de perguntas não relacionadas às disciplinas do mesmo conjunto de dados \parencite{zaman2024dataset}, a avaliação da qualidade das respostas e dos resumos gerados e a aplicação da plataforma em turmas reais, com perguntas em português e a participação de professores. Por fim, propõe-se a identificação da etapa em que surgem as divergências de classificação, a partir dos fragmentos recuperados registrados para cada pergunta.
}
\clearpage

% === REFERÊNCIAS
\sectionCentralizadaComSumario{REFERÊNCIAS}
\printbibliography[heading=none]
\clearpage

% === ELEMENTOS PÓS-TEXTUAIS
\glossario
\clearpage
%\apendice
%\clearpage
%\anexos
%\clearpage
%\indice

\end{document}
